% Template KLTN cho SV trường ĐHKHTN
% Nguồn tham khảo: https://www.overleaf.com/project/5f0e9a442545fc0001635d7f
% Last update: 13/11/2025

% Chú ý:
% - Template được sử dụng tốt nhất trên Overleaf.
% - Nếu sử dụng template trên các nền tảng offline (Vd: TexLive và Texmaker) cần lưu ý những điều sau:
%   1. Trước khi build, xóa hết các file được tạo ra trong quá trình build trước đó, và build theo thứ tự: BIB > PDF > PDF.
%   2. Nếu cập nhật tài liệu tham khảo, cũng cần build lại theo cách trên.


\documentclass[twoside,a4paper,14pt,openright]{extreport}
% Font size
\usepackage[fontsize=13pt]{scrextend}
% Font tiếng Việt
\usepackage[T5]{fontenc}
\usepackage[utf8]{inputenc}
\DeclareTextSymbolDefault{\DH}{T1}


% Tài liệu tham khảo
\usepackage[
	sorting=nty,
	backend=bibtex,
	defernumbers=true]{biblatex}
\usepackage[unicode]{hyperref} % Bookmark tiếng Việt
\addbibresource{References/references.bib}

\makeatletter
\def\blx@maxline{77}
\makeatother

% Chèn hình, các hình trong luận văn được để trong thư mục Images/
\usepackage{graphicx}
\graphicspath{ {Images/} }
\usepackage{caption}
\usepackage[export]{adjustbox}

% Chèn và định dạng mã nguồn
\usepackage{listings}
\usepackage{color}
\definecolor{codegreen}{rgb}{0,0.6,0}
\definecolor{codegray}{rgb}{0.5,0.5,0.5}
\definecolor{codepurple}{rgb}{0.58,0,0.82}
\definecolor{backcolour}{rgb}{0.95,0.95,0.92}
\lstdefinestyle{mystyle}{
    backgroundcolor=\color{backcolour},   
    commentstyle=\color{codegreen},
    keywordstyle=\color{magenta},
    numberstyle=\tiny\color{codegray},
    stringstyle=\color{codepurple},
    basicstyle=\footnotesize,
    breakatwhitespace=false,         
    breaklines=true,                 
    captionpos=b,                    
    keepspaces=true,                 
    numbers=left,                    
    numbersep=5pt,                  
    showspaces=false,                
    showstringspaces=false,
    showtabs=false,                  
    tabsize=2
}
\lstset{style=mystyle}

% Chèn và định dạng mã giả
\usepackage{amsmath}
\usepackage{algorithm}
\usepackage[noend]{algpseudocode}
\makeatletter
\def\BState{\State\hskip-\ALG@thistlm}
\makeatother
\makeatletter
\renewcommand{\ALG@name}{Thuật toán}
\renewcommand{\algorithmicrequire}{\textbf{Đầu vào:}}
\renewcommand{\algorithmicensure}{\textbf{Đầu ra:}}
\makeatother


% Bảng biểu
\usepackage{multirow}
\usepackage{array}
\newcolumntype{L}[1]{>{\raggedright\let\newline\\\arraybackslash\hspace{0pt}}m{#1}}
\newcolumntype{C}[1]{>{\centering\let\newline\\\arraybackslash\hspace{0pt}}m{#1}}
\newcolumntype{R}[1]{>{\raggedleft\let\newline\\\arraybackslash\hspace{0pt}}m{#1}}
\usepackage{diagbox}
\usepackage[table,xcdraw]{xcolor}
\usepackage{pmboxdraw}
\def\SFii{\textcolor{black}{\textSFii\textSFx}} %└
\def\SFviii{\textcolor{black}{\textSFviii\textSFx}} %├
\def\SFx{\textcolor{black}{\textSFx}} %─
\def\SFxi{\textcolor{black}{\textSFxi}} %|
 
% Đổi tên mặc định
\renewcommand{\chaptername}{CHƯƠNG}
\renewcommand{\figurename}{Hình}
\renewcommand{\tablename}{Bảng}
\renewcommand{\contentsname}{MỤC LỤC}
\renewcommand{\listfigurename}{DANH SÁCH CÁC HÌNH}
\renewcommand{\listtablename}{DANH SÁCH CÁC BẢNG}
\renewcommand{\appendixname}{PHỤ LỤC}

% Định dạng chapter
\usepackage{titlesec}

\setcounter{secnumdepth}{3}

\titleformat{\chapter}
    [block]
    {\filcenter\normalfont\bfseries\normalsize}{\chaptername \ \thechapter :}{10pt}{\normalsize}
\titlespacing*{\chapter}{0pt}{-10pt}{10pt} %khoảng cách giữa chapter và đầu trang

\titleformat{\section}
    {\normalfont\bfseries\normalsize}{\thesection}{1em}{}
 
\titleformat{\subsection}
    {\normalfont\bfseries\normalsize}{\thesubsection}{1em}{}

\titleformat{\subsubsection}
    {\normalfont\bfseries\normalsize}{\thesubsubsection}{1em}{}

% Dãn dòng 1.5
\usepackage{setspace}
\onehalfspacing

% Thụt vào đầu dòng
\usepackage{indentfirst}

% Bottom right page numbering
\usepackage{fancyhdr} % for use of \pageref{LastPage}
\pagestyle{fancy} % Turn on the style
\fancyhf{} % Start with clearing everything in the header and footer
% Set the right side of the footer to be the page number
\fancyfoot[R]{\thepage}
\renewcommand{\headrulewidth}{0pt}%
% Redefine plain style, which is used for titlepage and chapter beginnings
% From https://tex.stackexchange.com/a/30230/828
\fancypagestyle{plain}{%
    \renewcommand{\headrulewidth}{0pt}%
    \fancyhf{}%
    \fancyfoot[R]{\thepage}%
}

% Canh lề
\usepackage[
  twoside,
  top=20mm,
  bottom=20mm,
  left=25mm,
  right=20mm,
  footskip = 15mm,
  includefoot]{geometry}

% Trang bìa
\usepackage{tikz}
\usetikzlibrary{calc}
\newcommand\HRule{\rule{\textwidth}{1pt}}

% Danh sách chữ viết tắt
\usepackage{myacronyms}

% ifelse
\usepackage{ifthen}

\newcommand\myemptypage{
\newpage
\thispagestyle{empty}
\mbox{}
\newpage
}

% ========================================================================================= %
% CHÚ Ý: Thông tin chung về KLTN - sinh viên điền vào đây để tự động update các trang khác  %
% ========================================================================================= %
\newcommand{\tenSV}{Trịnh~Huy~Hoàng} % Dấu ~ là khoảng trắng không được tách
\newcommand{\mssv}{21207156}
\newcommand{\tenKL}{THIẾT~KẾ~VÀ~KIỂM~CHỨNG~VI~XỬ~LÝ~RISC-V~TÍCH~HỢP~TRÊN~CHIP~(SoC)~VỚI~GIAO~DIỆN~AXI-Lite~VÀ~AHB-Lite}
\newcommand{\tenGVHD}{(GS/PGS).(TS/ThS/CN). Nguyễn Văn A}
\newcommand{\tenBM}{Hệ thống nhúng và Vi mạch}

\begin{document}

% fix anchor bug
\hypersetup{pageanchor=false}
\begin{titlepage}

\begin{center}
\textbf{ĐẠI HỌC QUỐC GIA THÀNH PHỐ HỒ CHÍ MINH\\
TRƯỜNG ĐẠI HỌC KHOA HỌC TỰ NHIÊN\\
KHOA ĐIỆN TỬ - VIỄN THÔNG}\\[2cm]


{ \Large \bfseries \tenSV \\[2cm] } 

% Tên đề tài Khóa luận tốt nghiệp/Đồ án tốt nghiệp

{ \Large \bfseries \tenKL \\[5cm]} 


% Chọn trong các dòng sau
\large KHÓA LUẬN TỐT NGHIỆP CỬ NHÂN\\
%\large ĐỒ ÁN TỐT NGHIỆP CỬ NHÂN\\
%\large THỰC TẬP TỐT NGHIỆP CỬ NHÂN\\
\large NGÀNH KỸ THUẬT ĐIỆN TỬ - VIỄN THÔNG\\
\large CHUYÊN NGÀNH: \tenBM (Nếu là chương trình đề án thì bỏ dòng này)


\begin{tikzpicture}[remember picture, overlay]
  \draw[line width = 2pt] ($(current page.north west) + (2cm,-2cm)$) rectangle ($(current page.south east) + (-1.5cm,2cm)$);
\end{tikzpicture}

\vfill
Tp. Hồ Chí Minh - Năm 2025

\end{center}

\pagebreak

\myemptypage

\begin{center}
\textbf{ĐẠI HỌC QUỐC GIA THÀNH PHỐ HỒ CHÍ MINH\\
TRƯỜNG ĐẠI HỌC KHOA HỌC TỰ NHIÊN\\
KHOA ĐIỆN TỬ - VIỄN THÔNG}\\[2cm]

% Chọn trong các dòng sau
\large \textbf{KHÓA LUẬN TỐT NGHIỆP CỬ NHÂN}\\
% \large ĐỒ ÁN TỐT NGHIỆP CỬ NHÂN\\
\large \textbf{NGÀNH KỸ THUẬT ĐIỆN TỬ - VIỄN THÔNG}\\%[2cm] (bỏ comment nếu là chương trình đề án)
\large \textbf{CHUYÊN NGÀNH: \tenBM} (Nếu là chương trình đề án thì bỏ dòng này)\\[2cm]


% Tên đề tài Khóa luận tốt nghiệp/Đồ án tốt nghiệp
{ \Large \bfseries \tenKL \\[2cm] } 

{\bfseries Họ và tên sinh viên: \tenSV\\
Mã số SV: \mssv\\[2cm]}


NGƯỜI HƯỚNG DẪN KHOA HỌC\\
1. \tenGVHD \\
2. (Nếu 1 người thì bỏ đánh số)



% \begin{tikzpicture}[remember picture, overlay]
%   \draw[line width = 2pt] ($(current page.north west) + (2cm,-2cm)$) rectangle ($(current page.south east) + (-1.5cm,2cm)$);
% \end{tikzpicture}

\vfill
Tp. Hồ Chí Minh - Năm 2025

\end{center}
\thispagestyle{empty}

\myemptypage

\end{titlepage}
\hypersetup{pageanchor=true}
% Sasu trang Title, các bạn chèn nhận xét gủa GVHD và GVPB. Nhận xét sẽ được giáo vụ phát sau buổi bảo vệ để các bạn đóng quyển.

\pagenumbering{roman} % Đánh số i, ii, iii, ...

% \addcontentsline{toc}{chapter}{Lời cam đoan}
% \include{Appendix/reassurances}
\chapter*{XÁC NHẬN HOÀN TẤT CHỈNH SỬA BÁO CÁO TỐT NGHIỆP}

{\textbf{TÊN ĐỀ TÀI: TÊN KHÓA LUẬN TỐT NGHIỆP TIẾNG VIỆT}.}

{Tên tiếng anh: Tên khóa luận tiếng anh.}\\[.1 cm]

Họ và tên sinh viên: \tenSV

Mã số sinh viên: \mssv \\[.1 cm]

\textbf{Hội đồng đánh giá họp ngày..... tháng..... năm.....}

\begin{enumerate}
	\item Chủ tịch: 
	\item Thư ký: 
	\item Ủy viên: 
	\item Ủy viên: 
	\item Ủy viên: 
\end{enumerate}

\vspace{.2 cm}
Người hướng dẫn: \tenGVHD

Xác nhận của người hướng dẫn:

\chapter*{LỜI CẢM ƠN}
\label{thanks}
\addcontentsline{toc}{chapter}{LỜI CẢM ƠN}

(Cảm ơn giáo viên hướng dẫn).

(Cảm ơn gia đình).

(Cảm ơn giáo viên trong khoa).

(Cảm ơn bạn bè, người yêu, vv).

\vspace{1cm}
\begin{flushright}
\begin{tabular}{c}
Thành phố Hồ Chí Minh, ngày 13 tháng 11 năm 2025\\[0.5cm]
\textbf{Sinh viên thực hiện}\\[2cm]
\textbf{\tenSV}
\end{tabular}
\end{flushright}


\include{Appendix/reassurances}

\chapter*{TÓM TẮT}
\label{tomtat}
\addcontentsline{toc}{chapter}{TÓM TẮT}

Khóa luận này trình bày quá trình thiết kế và kiểm chứng một hệ thống trên chip (SoC) dựa trên kiến trúc tập lệnh RISC-V mở, nhằm mục tiêu nghiên cứu và học thuật. Hệ thống bao gồm lõi vi xử lý RV32I+Zicsr với kiến trúc pipeline 7 tầng hoạt động ở tần số 1\,GHz, tích hợp bộ nhớ SRAM nội bộ (IMEM 64\,KB và DMEM 64\,KB), giao thức bus AXI4-Lite (1\,GHz) và AHB-Lite (500\,MHz), cùng cơ chế đồng bộ hóa qua biên miền xung nhịp (CDC) bằng FIFO bất đồng bộ với con trỏ Gray code. Hệ thống ngắt bao gồm PLIC 6 nguồn tương thích chuẩn SiFive và khối Zicsr xử lý ngắt/ngoại lệ M-mode. Toàn bộ RTL được viết bằng SystemVerilog có thể tổng hợp và mô phỏng bằng Icarus Verilog.

Quá trình kiểm chứng được tổ chức thành nhiều pha tăng dần về độ phức tạp: từ unit test (hơn 430 test case), integration test (163 test case), system test (37 chương trình), đến bộ kiểm thử tuân thủ chính thức riscv-arch-test. Kết quả: tất cả test case PASS và đạt 37/38 PASS trong bộ kiểm thử tuân thủ RV32I (một bài SKIP do giới hạn kích thước code của framework). Trong quá trình kiểm thử, năm bug RTL quan trọng được phát hiện và sửa chữa, bao gồm gap-4 RAW hazard, IMEM stall mismatch, CSR-use hazard stall, ghost instruction sau flush, và bus stall kết hợp load-use hazard.

\textbf{Từ khóa} --- RISC-V, SoC, pipeline, AXI-Lite, AHB-Lite, CDC, PLIC, SystemVerilog, kiểm chứng RTL.

\begin{center}
    \textbf{ABSTRACT}
\end{center}

This thesis presents the design and verification of a RISC-V-based System-on-Chip (SoC) for academic research purposes. The system features an RV32I+Zicsr processor with a 7-stage in-order pipeline targeting 1\,GHz, integrated with on-chip SRAM (64\,KB IMEM and 64\,KB DMEM), AXI4-Lite (1\,GHz) and AHB-Lite (500\,MHz) bus interfaces, and safe Clock Domain Crossing (CDC) via asynchronous FIFO with Gray-code pointers. The interrupt subsystem includes a SiFive-compatible 6-source PLIC and a Zicsr block handling M-mode exceptions and interrupts. All RTL is written in synthesizable SystemVerilog and simulated using Icarus Verilog.

Verification follows a multi-level strategy: unit tests (430+ test cases), integration tests (163 test cases), system-level programs (37 programs), and the official riscv-arch-test compliance suite. All test cases pass, and 37/38 RV32I compliance tests pass (one SKIP due to framework code-size limitations). Five RTL bugs were discovered and fixed during testing, including gap-4 RAW hazard, IMEM stall mismatch, CSR-use hazard stall, ghost instruction after flush, and a concurrent bus-stall/load-use hazard interaction.

\textbf{Keywords} --- RISC-V, SoC, pipeline, AXI-Lite, AHB-Lite, CDC, PLIC, SystemVerilog, RTL verification.


% \phantomsection
% \addcontentsline{toc}{chapter}{Đề cương chi tiết}
% \include{Appendix/decuong}

% Mục lục, danh sách hình, danh sách bảng
\cleardoublepage
\phantomsection
% \addcontentsline{toc}{chapter}{MỤC LỤC}
\tableofcontents
% Danh sách chữ viết tắt
\printglossary[type=\acronymtype,
                title=DANH SÁCH CHỮ VIẾT TẮT, 
                toctitle=DANH SÁCH CHỮ VIẾT TẮT,
                nonumberlist
                ]

\cleardoublepage
\newpage
\phantomsection
\addcontentsline{toc}{chapter}{\listfigurename}
\listoffigures

\cleardoublepage
\newpage
\phantomsection
\addcontentsline{toc}{chapter}{\listtablename}
\listoftables

\cleardoublepage

\pagenumbering{arabic} % Đánh số 1, 2, 3, ...

% Các chương nội dung
\chapter{GIỚI THIỆU}
\label{Chapter1}

\section{Bối Cảnh và Động Lực Nghiên Cứu}

Trong bối cảnh ngành công nghiệp vi điện tử đang phát triển nhanh chóng, xu hướng thiết kế phần cứng mở (open-source hardware) ngày càng được quan tâm rộng rãi. Kiến trúc tập lệnh RISC-V (Reduced Instruction Set Computer -- Version 5), ra đời tại Đại học California, Berkeley vào năm 2010, đã trở thành một trong những nền tảng quan trọng nhất trong xu hướng này \cite{riscv-spec}. Khác với các kiến trúc thương mại như ARM hay x86, RISC-V hoàn toàn mở, không thu phí bản quyền và được quản lý bởi tổ chức phi lợi nhuận RISC-V International. Điều này tạo điều kiện thuận lợi để các cơ sở học thuật và doanh nghiệp có thể tự do nghiên cứu, triển khai và tùy chỉnh kiến trúc theo nhu cầu cụ thể.

Theo báo cáo thị trường của Semico Research, số lõi vi xử lý RISC-V được tích hợp vào sản phẩm thương mại dự kiến đạt 62 tỷ đơn vị vào năm 2025 \cite{semico-riscv-growth}. Sự phát triển vượt bậc này được thúc đẩy bởi nhu cầu ngày càng cao về các hệ thống nhúng hiệu suất cao, tiêu thụ điện năng thấp, và có thể tùy chỉnh linh hoạt -- từ các thiết bị IoT nhỏ gọn đến các hệ thống xử lý dữ liệu phức tạp trong trung tâm dữ liệu. Trong bối cảnh đó, việc nghiên cứu và thiết kế một hệ thống trên chip (\acrfull{soc}) hoàn chỉnh dựa trên RISC-V có giá trị học thuật và thực tiễn cao.

Thiết kế một SoC không chỉ đòi hỏi việc xây dựng lõi vi xử lý mà còn bao gồm tích hợp các giao thức bus công nghiệp chuẩn để kết nối với ngoại vi, xử lý ngắt và ngoại lệ, đồng bộ hóa giữa các miền xung nhịp khác nhau, và kiểm chứng toàn diện để đảm bảo tính đúng đắn của thiết kế. Đây là những kỹ năng thiết yếu trong quy trình thiết kế vi mạch tích hợp (\acrshort{asic}/FPGA) hiện đại.

\section{Mục Tiêu và Phạm Vi}

Khóa luận này đặt ra mục tiêu thiết kế và kiểm chứng một hệ thống SoC hoàn chỉnh dựa trên kiến trúc RISC-V, bao gồm:

\begin{enumerate}
    \item \textbf{Lõi CPU:} Xây dựng vi xử lý RV32I với phần mở rộng Zicsr theo đặc tả RISC-V, tổ chức theo kiến trúc pipeline 7 tầng với tần số mục tiêu 1\,GHz.
    \item \textbf{Tích hợp bus chuẩn công nghiệp:} Kết nối CPU với ngoại vi thông qua hai giao thức bus -- AXI4-Lite (AMBA~4.0) và AHB-Lite (AMBA~3.0) -- đại diện cho hai tiêu chuẩn phổ biến nhất trong thiết kế SoC hiện đại \cite{amba-axi-spec, amba-ahb-spec}.
    \item \textbf{Xử lý đa miền xung nhịp:} Triển khai cơ chế Clock Domain Crossing (\acrshort{cdc}) sử dụng FIFO bất đồng bộ với con trỏ Gray code để truyền dữ liệu an toàn giữa miền 1\,GHz và 500\,MHz.
    \item \textbf{Hệ thống ngắt và ngoại lệ:} Triển khai \acrfull{plic} tương thích chuẩn SiFive và khối Zicsr xử lý ngắt/ngoại lệ M-mode theo đặc tả RISC-V Privileged Architecture \cite{riscv-priv-spec}.
    \item \textbf{Kiểm chứng toàn diện:} Xây dựng môi trường kiểm thử nhiều lớp và đạt chuẩn tuân thủ bộ kiểm thử chính thức RV32I của tổ chức RISC-V International \cite{riscv-arch-test}.
\end{enumerate}

Phạm vi thiết kế giới hạn ở mức RTL (Register Transfer Level) sử dụng ngôn ngữ SystemVerilog có thể tổng hợp. Toàn bộ mô phỏng được thực hiện bằng công cụ Icarus Verilog~12 \cite{icarus-verilog}.

\section{Phương Pháp Thực Hiện}

Quy trình thiết kế và kiểm chứng trong khóa luận này theo mô hình phát triển từ trên xuống (top-down) kết hợp kiểm chứng từ dưới lên (bottom-up verification):

\begin{itemize}
    \item \textbf{Pha 1 -- Đặc tả:} Xác định kiến trúc hệ thống, bản đồ bộ nhớ, giao thức bus và các yêu cầu chức năng chi tiết.
    \item \textbf{Pha 2 -- Thiết kế RTL:} Viết các module SystemVerilog theo từng khối chức năng, tuân thủ các quy tắc coding cho synthesizable design và yêu cầu của công cụ mô phỏng.
    \item \textbf{Pha 3 -- Kiểm chứng phân cấp:} Bắt đầu từ unit test từng module, tiếp theo là integration test các subsystem, và system test toàn bộ SoC với nhiều chương trình thử nghiệm khác nhau.
    \item \textbf{Pha 4 -- Kiểm thử tuân thủ:} Sử dụng bộ kiểm thử chính thức riscv-arch-test để xác nhận tính tuân thủ đặc tả RV32I.
\end{itemize}

\section{Các Công Cụ Sử Dụng}

Bảng~\ref{tab:tools} liệt kê các công cụ phần mềm chính được sử dụng trong quá trình thực hiện khóa luận.

\begin{table}[htbp]
\caption{Công cụ phần mềm sử dụng trong dự án}
\label{tab:tools}
\begin{center}
\begin{tabular}{|l|l|l|}
\hline
\textbf{Công cụ} & \textbf{Phiên bản} & \textbf{Mục đích} \\ \hline
Icarus Verilog & 12 & Mô phỏng RTL SystemVerilog \\ \hline
riscv64-unknown-elf-gcc & 13.2.0 & Biên dịch chương trình RISC-V \\ \hline
GTKWave & 3.3.x & Xem dạng sóng (waveform) VCD \\ \hline
riscv-arch-test & 2.x (old-framework) & Bộ kiểm thử tuân thủ RV32I \\ \hline
GNU Make & 4.x & Tự động hóa quy trình build/test \\ \hline
\end{tabular}
\end{center}
\end{table}

\section{Cấu Trúc Khóa Luận}

Nội dung khóa luận được tổ chức thành sáu chương, trình bày tuần tự từ bối cảnh nghiên cứu, tổng quan lý thuyết, thiết kế, triển khai thực nghiệm đến phân tích, đánh giá kết quả và kết luận. Cụ thể:

\begin{itemize}
    \item \textbf{Chương 1 -- Giới thiệu:} Trình bày bối cảnh nghiên cứu, động lực thực hiện đề tài, mục tiêu, phạm vi, công cụ sử dụng và tóm lược cấu trúc toàn khóa luận.
    \item \textbf{Chương 2 -- Tổng quan nghiên cứu:} Điểm qua các dự án RISC-V SoC tiêu biểu trong và ngoài nước, phân tích điểm tương đồng và khác biệt so với thiết kế đề xuất.
    \item \textbf{Chương 3 -- Cơ sở lý thuyết:} Trình bày nền tảng lý thuyết về kiến trúc tập lệnh RISC-V, kỹ thuật pipeline, các giao thức bus AXI-Lite và AHB-Lite, cơ chế CDC và xử lý ngắt.
    \item \textbf{Chương 4 -- Thiết kế và triển khai hệ thống:} Mô tả chi tiết kiến trúc và cách triển khai từng thành phần của SoC, bao gồm pipeline CPU, bus interface, CDC, PLIC và Zicsr.
    \item \textbf{Chương 5 -- Kết quả thực nghiệm:} Trình bày môi trường kiểm thử, phân tích kết quả từng pha kiểm thử, phân tích dạng sóng minh họa và đánh giá tổng thể hiệu quả thiết kế.
    \item \textbf{Chương 6 -- Kết luận:} Tóm tắt kết quả đạt được, nêu hạn chế và đề xuất hướng phát triển tiếp theo.
\end{itemize}

Bên cạnh đó, khóa luận còn đính kèm phần \textbf{Tài liệu tham khảo} ở cuối nhằm cung cấp thông tin chi tiết cho việc tra cứu.


\chapter{TỔNG QUAN NGHIÊN CỨU}
\label{Chapter2}

\section{Các Dự Án RISC-V SoC Tiêu Biểu}

Kể từ khi đặc tả RISC-V được công bố mở năm 2014, một số lượng lớn các dự án thiết kế SoC dựa trên RISC-V đã xuất hiện từ các tổ chức học thuật, công ty công nghiệp và cộng đồng mã nguồn mở. Việc khảo sát các dự án này cung cấp nền tảng để định vị thiết kế của khóa luận trong bức tranh tổng thể của lĩnh vực.

\subsection{SiFive FE310 và HiFive1}

SiFive, công ty khởi nghiệp được sáng lập bởi các nhà nghiên cứu RISC-V từ UC Berkeley, đã phát triển dòng vi xử lý Freedom E310 (FE310) -- một trong những SoC thương mại đầu tiên dựa trên RISC-V \cite{sifive-fe310}. FE310 sử dụng lõi E31 với pipeline 5 tầng, kiến trúc RV32IMAC, và được tích hợp trên board phát triển HiFive1. Đặc điểm nổi bật của FE310 bao gồm PLIC (Platform-Level Interrupt Controller) tương thích đặc tả SiFive, CLINT (Core Local Interruptor), và các ngoại vi SPI/UART/GPIO. Tuy nhiên, FE310 sử dụng giao thức TileLink thay vì AMBA AXI/AHB, nên không trực tiếp tương thích với hệ sinh thái ngoại vi AMBA phổ biến.

\subsection{Dự Án OpenTitan}

OpenTitan là dự án nguồn mở của Google và các đối tác công nghiệp, nhằm tạo ra một Root of Trust chip theo thiết kế mở \cite{opentitan}. OpenTitan sử dụng lõi Ibex (dựa trên RI5CY) và giao thức bus TileLink-UL cho kết nối nội bộ. Điểm đáng chú ý của OpenTitan là hệ thống quản lý thanh ghi ngoại vi (SFR -- Special Function Register) được chuẩn hóa theo OpenTitan Register Tool, cung cấp giao diện nhất quán cho mọi ngoại vi. Thiết kế đề xuất trong khóa luận này lấy cảm hứng từ cách tổ chức thanh ghi của OpenTitan để xây dựng SFR Standard Register Map.

\subsection{VexRiscv}

VexRiscv là một vi xử lý RISC-V linh hoạt được viết bằng SpinalHDL, được thiết kế theo kiến trúc plugin cho phép tùy chỉnh sâu pipeline và tính năng \cite{vexriscv}. VexRiscv hỗ trợ từ pipeline 2 tầng đơn giản đến pipeline 5 tầng đầy đủ với cache L1 và MMU. Giao thức bus sử dụng WishBone hoặc AXI4. Điểm khác biệt so với thiết kế trong khóa luận là VexRiscv không hỗ trợ AHB-Lite natively và không tích hợp CDC cho đa miền xung nhịp.

\subsection{PicoRV32}

PicoRV32 là một lõi RISC-V RV32IMC nhỏ gọn, được viết thuần túy bằng Verilog và tối ưu cho FPGA \cite{picorv32}. PicoRV32 sử dụng kiến trúc không pipeline (hay pipeline tối giản), ưu tiên tính đơn giản và tài nguyên phần cứng thấp. Giao thức bus riêng của nó (PicoRV32 Memory Interface) không tương thích với AMBA. Thiết kế này phù hợp cho các ứng dụng nhúng cực nhỏ nhưng không phù hợp cho các hệ thống yêu cầu hiệu suất cao và kết nối ngoại vi phong phú.

\subsection{Các Dự Án Học Thuật Khác}

Trong giới học thuật, nhiều trường đại học đã công bố thiết kế RISC-V SoC như một phần của nghiên cứu. Điển hình là BOOM (Berkeley Out-of-Order Machine) từ UC Berkeley -- một vi xử lý RISC-V 64-bit out-of-order phức tạp phục vụ nghiên cứu về vi kiến trúc hiệu suất cao \cite{boom}. Tuy nhiên, BOOM tập trung vào nghiên cứu kiến trúc thuần túy và không đặt trọng tâm vào tích hợp bus chuẩn công nghiệp hay kiểm chứng đa miền xung nhịp. Bài báo~\cite{riscv-survey} cung cấp một khảo sát toàn diện về các triển khai và mở rộng RISC-V tính đến năm 2023.

\section{So Sánh Với Thiết Kế Đề Xuất}

Bảng~\ref{tab:comparison} so sánh đặc điểm thiết kế của các dự án tiêu biểu với thiết kế được đề xuất trong khóa luận này.

\begin{table}[htbp]
\caption{So sánh các dự án RISC-V SoC tiêu biểu}
\label{tab:comparison}
\begin{center}
\begin{tabular}{|l|c|c|c|c|c|}
\hline
\textbf{Đặc điểm} & \textbf{FE310} & \textbf{VexRiscv} & \textbf{PicoRV32} & \textbf{OpenTitan} & \textbf{Đề xuất} \\ \hline
ISA & RV32IMAC & RV32I/IM(A) & RV32IMC & RV32I & RV32I+Zicsr \\ \hline
Pipeline & 5 tầng & 2--5 tầng & Không & 2 tầng & \textbf{7 tầng} \\ \hline
Bus AXI-Lite & Không & Có & Không & Không & \textbf{Có} \\ \hline
Bus AHB-Lite & Không & Không & Không & Không & \textbf{Có} \\ \hline
CDC đa miền & Không & Không & Không & Không & \textbf{Có (FIFO)} \\ \hline
PLIC tích hợp & Có & Tùy chọn & Không & Có & \textbf{Có (6 nguồn)} \\ \hline
SFR chuẩn hóa & Không & Không & Không & Có & \textbf{Có} \\ \hline
RV32I compliance & Không rõ & Có & Có & Có & \textbf{37/38 PASS} \\ \hline
RTL Language & Chisel & SpinalHDL & Verilog & SV & \textbf{SystemVerilog} \\ \hline
\end{tabular}
\end{center}
\end{table}

Từ bảng so sánh, có thể thấy thiết kế đề xuất có một số điểm khác biệt đáng chú ý so với các giải pháp hiện có:

\begin{itemize}
    \item \textbf{Pipeline 7 tầng:} Hầu hết các thiết kế học thuật và nhúng chọn pipeline 2--5 tầng để đơn giản hóa xử lý hazard. Pipeline 7 tầng của thiết kế này tăng throughput tiềm năng và là thách thức lớn hơn về mặt xử lý RAW hazard, forwarding, và precise exception.
    \item \textbf{Tích hợp đồng thời AXI-Lite và AHB-Lite:} Không có thiết kế nào trong danh sách khảo sát hỗ trợ cả hai giao thức AMBA trong cùng một SoC với CDC. Đây là tính năng đặc trưng của thiết kế này, phản ánh thực tế công nghiệp nơi các SoC phải kết nối với nhiều loại ngoại vi từ các nhà cung cấp khác nhau.
    \item \textbf{CDC với Async FIFO:} Việc triển khai CDC bằng FIFO bất đồng bộ Gray code giữa hai miền xung nhịp là một giải pháp kỹ thuật phức tạp, phổ biến trong ASIC nhưng hiếm thấy trong các dự án học thuật cấp đại học cử nhân.
    \item \textbf{SFR Standard Register Map:} Thiết kế một chuẩn thanh ghi thống nhất cho tất cả ngoại vi (lấy cảm hứng từ OpenTitan) giúp đảm bảo tính mở rộng của hệ thống mà không cần thay đổi RTL lõi.
\end{itemize}

\section{Khoảng Trống Nghiên Cứu}

Qua khảo sát tài liệu, có thể xác định một số khoảng trống mà thiết kế đề xuất nhằm lấp đầy trong phạm vi một khóa luận cử nhân:

Thứ nhất, đa số các thiết kế học thuật RISC-V tập trung vào một giao thức bus duy nhất (hoặc giao thức bus riêng). Việc tích hợp đồng thời AXI-Lite (1\,GHz) và AHB-Lite (500\,MHz) trong cùng một SoC với CDC an toàn là một kịch bản kỹ thuật thực tế nhưng chưa được trình bày rõ ràng ở cấp độ học thuật với đầy đủ kiểm chứng.

Thứ hai, kiểm chứng đa cấp độ -- từ unit test qua integration test đến system test và compliance test chính thức -- thường không được trình bày đầy đủ trong các công bố học thuật về thiết kế vi xử lý. Khóa luận này đặt trọng tâm vào quy trình kiểm chứng có hệ thống như một đóng góp phương pháp luận.


\chapter{CƠ SỞ LÝ THUYẾT}
\label{Chapter3}

\section{Kiến Trúc Tập Lệnh RISC-V}

\subsection{Tổng Quan RISC-V ISA}

RISC-V là một kiến trúc tập lệnh mở (open ISA), được thiết kế theo triết lý RISC (Reduced Instruction Set Computer): số lượng lệnh nhỏ, định dạng cố định, và thực thi trong một chu kỳ trên pipeline lý tưởng \cite{riscv-spec}. Điểm mạnh của RISC-V là tính module: một tập lệnh cơ sở bắt buộc (RV32I hoặc RV64I) kết hợp với các phần mở rộng tùy chọn (M, A, F, D, C, Zicsr, v.v.) cho phép thiết kế tùy chỉnh linh hoạt theo nhu cầu ứng dụng.

Phần mở rộng \textbf{RV32I} (32-bit Integer Base) bao gồm 47 lệnh cơ bản, đủ để chạy phần lớn phần mềm ứng dụng. RV32I định nghĩa:

\begin{itemize}
    \item 32 thanh ghi nguyên x0--x31, mỗi thanh ghi 32 bit. Thanh ghi x0 luôn bằng 0 (hardwired zero).
    \item Chiều rộng lệnh 32 bit cố định.
    \item Không gian địa chỉ 32 bit, byte-addressable, word-aligned instruction fetch.
\end{itemize}

\subsection{Định Dạng Lệnh RV32I}

RISC-V định nghĩa sáu định dạng mã hóa lệnh, tất cả đều 32 bit. Hình~\ref{fig:riscv-formats} minh họa cấu trúc của từng định dạng.

\begin{figure}[htp]
\centering
\fbox{\parbox{0.85\textwidth}{\centering\textit{[Chèn hình: Sáu định dạng mã hóa lệnh RV32I -- R, I, S, B, U, J-type với các trường opcode, funct3, funct7, rs1, rs2, rd, imm]}\vspace{2.5cm}}}
\caption{Sáu định dạng mã hóa lệnh RV32I}
\label{fig:riscv-formats}
\end{figure}

Bảng~\ref{tab:instr-groups} tóm tắt các nhóm lệnh RV32I.

\begin{table}[htbp]
\caption{Các nhóm lệnh RV32I}
\label{tab:instr-groups}
\begin{center}
\begin{tabular}{|l|l|l|}
\hline
\textbf{Nhóm} & \textbf{Lệnh tiêu biểu} & \textbf{Chức năng} \\ \hline
Số học & ADD, SUB, ADDI & Cộng/trừ thanh ghi và immediate \\ \hline
Logic & AND, OR, XOR, ANDI, ORI, XORI & Phép logic bit \\ \hline
Dịch bit & SLL, SRL, SRA, SLLI, SRLI, SRAI & Dịch trái/phải logic/số học \\ \hline
So sánh & SLT, SLTU, SLTI, SLTIU & Ghi kết quả 0/1 vào rd \\ \hline
Load & LW, LH, LB, LHU, LBU & Đọc từ bộ nhớ \\ \hline
Store & SW, SH, SB & Ghi vào bộ nhớ \\ \hline
Nhánh & BEQ, BNE, BLT, BGE, BLTU, BGEU & Rẽ nhánh có điều kiện \\ \hline
Nhảy & JAL, JALR & Nhảy không điều kiện \\ \hline
PC-relative & AUIPC & Cộng imm dịch 12 bit vào PC \\ \hline
Immediate & LUI & Nạp 20-bit immediate vào rd[31:12] \\ \hline
Hệ thống & ECALL, EBREAK & Gọi hệ điều hành / gỡ lỗi \\ \hline
\end{tabular}
\end{center}
\end{table}

\subsection{Phần Mở Rộng Zicsr}

Phần mở rộng Zicsr thêm sáu lệnh đọc/ghi các thanh ghi điều khiển-trạng thái (CSR -- Control and Status Register) và định nghĩa cơ chế xử lý ngắt/ngoại lệ ở M-mode (Machine mode) \cite{riscv-priv-spec}. Khi xảy ra trap (ngắt hoặc ngoại lệ), phần cứng tự động thực hiện chuỗi hành động:

\begin{enumerate}
    \item Lưu PC hiện tại vào \texttt{mepc}.
    \item Ghi nguyên nhân vào \texttt{mcause} (bit~31 phân biệt interrupt/exception).
    \item Lưu trạng thái ngắt \texttt{MIE} vào \texttt{MPIE} trong \texttt{mstatus}.
    \item Tắt ngắt toàn cục (\texttt{MIE} $\leftarrow$ 0).
    \item Chuyển PC đến vector trap (\texttt{mtvec + 4×cause} trong vectored mode).
\end{enumerate}

Lệnh MRET phục hồi trạng thái trước trap: PC $\leftarrow$ \texttt{mepc}, \texttt{MIE} $\leftarrow$ \texttt{MPIE}.

\section{Kiến Trúc Pipeline Processor}

\subsection{Nguyên Lý Pipeline}

Pipeline là kỹ thuật tổ chức CPU theo nhiều tầng (stage), mỗi tầng thực hiện một phần của chu trình thực thi. Các tầng hoạt động song song, xử lý các lệnh khác nhau cùng một lúc (instruction-level parallelism). Throughput lý tưởng đạt 1 lệnh/chu kỳ (1 IPC).

\begin{figure}[htp]
\centering
\fbox{\parbox{0.85\textwidth}{\centering\textit{[Chèn hình: Minh họa pipeline 7 tầng -- các lệnh I1, I2, I3 đi qua các tầng IF1, IF2, ID, EX, MEM1, MEM2, WB theo thời gian, thể hiện instruction-level parallelism]}\vspace{3cm}}}
\caption{Nguyên lý hoạt động pipeline 7 tầng}
\label{fig:pipeline-timing}
\end{figure}

\subsection{Vấn Đề Hazard và Giải Pháp}

Trong pipeline, \textbf{hazard} là điều kiện ngăn lệnh tiếp theo thực thi ngay ở chu kỳ tiếp theo. Có ba loại hazard chính:

\textbf{1. Data hazard (RAW -- Read After Write):} Lệnh sau cần đọc dữ liệu mà lệnh trước chưa ghi xong. Giải pháp chính là \textbf{data forwarding} (bypass): dữ liệu được chuyển thẳng từ đầu ra tầng trước về đầu vào tầng EX mà không cần chờ ghi vào register file. Trường hợp đặc biệt là \textbf{load-use hazard}: lệnh LOAD có kết quả chỉ có ở tầng MEM2, nên lệnh ngay sau không thể forward kịp từ EX -- cần stall 1 chu kỳ.

\textbf{2. Structural hazard:} Hai lệnh cùng cần sử dụng cùng một tài nguyên phần cứng. Trong thiết kế pipeline với memory riêng biệt cho instruction và data (Harvard architecture), structural hazard hầu như không xuất hiện.

\textbf{3. Control hazard:} Branch/Jump -- khi CPU chưa biết địa chỉ đích, nó đã nạp lệnh sai vào pipeline. Giải pháp đơn giản là \textbf{flush} (xóa) các lệnh sai đã vào pipeline khi biết địa chỉ đích.

\section{Giao Thức AXI4-Lite}

\subsection{Tổng Quan}

AXI4-Lite (Advanced eXtensible Interface 4, Lite version) là giao thức bus của ARM, thuộc tiêu chuẩn AMBA~4.0 \cite{amba-axi-spec}. AXI4-Lite được thiết kế cho giao tiếp đơn giản với các thanh ghi điều khiển (control/status registers) và là tiêu chuẩn phổ biến nhất trong hệ sinh thái ARM và Xilinx/Intel FPGA.

\subsection{Cấu Trúc Kênh}

AXI4-Lite sử dụng năm kênh giao tiếp độc lập:

\begin{table}[htbp]
\caption{Năm kênh giao tiếp AXI4-Lite}
\label{tab:axi-channels}
\begin{center}
\begin{tabular}{|l|l|l|l|}
\hline
\textbf{Kênh} & \textbf{Hướng} & \textbf{Tín hiệu chính} & \textbf{Mô tả} \\ \hline
AW & Master$\rightarrow$Slave & AWADDR, AWVALID, AWREADY & Write Address \\ \hline
W & Master$\rightarrow$Slave & WDATA, WSTRB, WVALID, WREADY & Write Data \\ \hline
B & Slave$\rightarrow$Master & BRESP, BVALID, BREADY & Write Response \\ \hline
AR & Master$\rightarrow$Slave & ARADDR, ARVALID, ARREADY & Read Address \\ \hline
R & Slave$\rightarrow$Master & RDATA, RRESP, RVALID, RREADY & Read Data \\ \hline
\end{tabular}
\end{center}
\end{table}

\subsection{Giao Thức Handshake}

Mỗi kênh AXI sử dụng cơ chế handshake VALID/READY: giao dịch xảy ra khi cả VALID và READY đều bằng 1 tại cùng cạnh xung nhịp. Hình~\ref{fig:axi-waveform} minh họa timing của một giao dịch ghi và đọc AXI-Lite điển hình.

\begin{figure}[htp]
\centering
\fbox{\parbox{0.85\textwidth}{\centering\textit{[Chèn hình: Waveform AXI-Lite -- giao dịch ghi (AWVALID/AWREADY, WVALID/WREADY, BVALID/BREADY) và giao dịch đọc (ARVALID/ARREADY, RVALID/RREADY/RDATA) trên nền xung clk 1GHz]}\vspace{3.5cm}}}
\caption{Dạng sóng giao dịch AXI4-Lite ghi và đọc}
\label{fig:axi-waveform}
\end{figure}

\section{Giao Thức AHB-Lite}

\subsection{Tổng Quan}

AHB-Lite (Advanced High-performance Bus Lite) là phiên bản đơn master của AHB trong tiêu chuẩn AMBA~3.0 \cite{amba-ahb-spec}. Khác với AXI, AHB-Lite sử dụng một bus chia sẻ (shared bus) cho tất cả slave và cơ chế pipelining giữa address phase và data phase.

\subsection{Hai Phase Giao Dịch}

AHB-Lite phân chia một giao dịch thành hai phase kế tiếp nhau:

\begin{itemize}
    \item \textbf{Address Phase (Chu kỳ N):} Master phát địa chỉ (HADDR), loại giao dịch (HTRANS=NONSEQ), kích thước (HSIZE) và hướng (HWRITE) lên bus.
    \item \textbf{Data Phase (Chu kỳ N+1):} Master phát dữ liệu ghi (HWDATA); slave phát dữ liệu đọc (HRDATA) và tín hiệu ready (HREADY). Slave có thể chèn wait state bằng cách giữ HREADY=0.
\end{itemize}

Điểm đặc biệt của AHB-Lite là address phase của giao dịch kế tiếp có thể diễn ra đồng thời với data phase của giao dịch hiện tại, tạo ra pipelining hiệu quả. Hình~\ref{fig:ahb-waveform} minh họa timing AHB-Lite.

\begin{figure}[htp]
\centering
\fbox{\parbox{0.85\textwidth}{\centering\textit{[Chèn hình: Waveform AHB-Lite -- giao dịch ghi với hai phase riêng biệt: Address Phase (HADDR, HTRANS, HWRITE, HSEL) và Data Phase (HWDATA, HREADYOUT, HRESP)]}\vspace{3.5cm}}}
\caption{Dạng sóng giao dịch AHB-Lite}
\label{fig:ahb-waveform}
\end{figure}

\section{Clock Domain Crossing}

\subsection{Vấn Đề Metastability}

Khi hai miền xung nhịp (clock domain) có tần số hoặc pha khác nhau cần giao tiếp, tín hiệu từ miền nguồn có thể vi phạm yêu cầu setup/hold time của flip-flop trong miền đích. Khi đó flip-flop có thể rơi vào trạng thái \textbf{metastability} -- đầu ra không xác định, dao động trong khoảng thời gian không biết trước. Nếu metastability truyền sang các logic tầng sau, có thể gây lỗi hệ thống \cite{cummings-cdc}.

\subsection{2-FF Synchronizer}

Giải pháp cơ bản nhất cho CDC một bit là sử dụng \textbf{2-FF synchronizer}: tín hiệu bất đồng bộ đi qua hai flip-flop tuần tự, cả hai dùng xung nhịp của miền đích. Xác suất metastability giảm theo hàm mũ qua mỗi tầng FF bổ sung. Trong thiết kế này, 2-FF synchronizer được dùng cho tín hiệu ngắt từ AHB domain (500\,MHz) sang CPU domain (1\,GHz).

\subsection{Asynchronous FIFO với Gray Code}

Để truyền dữ liệu nhiều bit giữa hai miền xung nhịp, giải pháp an toàn là sử dụng \textbf{asynchronous FIFO} với con trỏ Gray code \cite{cummings-cdc}. Gray code đảm bảo chỉ một bit thay đổi mỗi bước đếm, nên con trỏ có thể được đồng bộ hóa an toàn qua 2-FF synchronizer mà không bị hiện tượng tear (đọc sai do nhiều bit thay đổi cùng lúc).

Hình~\ref{fig:async-fifo} minh họa kiến trúc và nguyên lý hoạt động của asynchronous FIFO với Gray code pointer.

\begin{figure}[htp]
\centering
\fbox{\parbox{0.85\textwidth}{\centering\textit{[Chèn hình: Sơ đồ khối Async FIFO -- bộ nhớ FIFO, write pointer (wr\_clk domain, mã Gray), read pointer (rd\_clk domain, mã Gray), hai 2-FF synchronizer đồng bộ con trỏ giữa hai miền, tín hiệu empty/full từ so sánh con trỏ]}\vspace{3.5cm}}}
\caption{Kiến trúc Asynchronous FIFO với Gray code pointer}
\label{fig:async-fifo}
\end{figure}

\section{Bộ Điều Khiển Ngắt Ưu Tiên (PLIC)}

\acrfull{plic} là một tiêu chuẩn kiến trúc ngắt của RISC-V cho phép quản lý tập trung nhiều nguồn ngắt ngoại vi trước khi đưa vào CPU \cite{sifive-plic}. PLIC hỗ trợ:

\begin{itemize}
    \item Nhiều nguồn ngắt với mức ưu tiên khác nhau (priority per source).
    \item Ngưỡng ngắt (threshold) cho phép CPU lọc các ngắt có ưu tiên thấp.
    \item Cơ chế claim/complete: CPU đọc CLAIM register để nhận ID ngắt, sau đó ghi vào COMPLETE khi xử lý xong.
\end{itemize}

PLIC đóng vai trò arbitrator trung gian: nhiều ngoại vi có thể kích hoạt ngắt đồng thời, nhưng PLIC chỉ chuyển ngắt có ưu tiên cao nhất (cao hơn ngưỡng) đến CPU qua tín hiệu \texttt{meip\_in}. Điều này giúp giảm tải cho CPU và chuẩn hóa cách xử lý ngắt trong hệ thống đa ngoại vi.


\chapter{THIẾT KẾ VÀ TRIỂN KHAI HỆ THỐNG}
\label{Chapter4}

\section{Kiến Trúc Tổng Thể SoC}

\subsection{Sơ Đồ Khối Hệ Thống}

Hệ thống SoC được thiết kế xoay quanh module tổng hợp \texttt{soc\_top}, bao gồm ba phân vùng chức năng chính: \textbf{CPU domain} (1\,GHz), \textbf{AXI-Lite subsystem} (1\,GHz), và \textbf{AHB-Lite subsystem} (500\,MHz với CDC). Module \texttt{soc\_top} không chứa bất kỳ logic datapath nào -- tất cả logic được đóng gói trong các module con và \texttt{soc\_top} chỉ kết dây (wire-up) giữa chúng. Triết lý thiết kế này giúp mỗi khối có thể kiểm thử độc lập.

\begin{figure}[htp]
\centering
\fbox{\parbox{0.85\textwidth}{\centering\textit{[Chèn hình: Sơ đồ khối tổng thể SoC -- CPU Pipeline (IF1, IF2, ID, EX, MEM1, MEM2, WB) + Hazard Unit + Forwarding + Zicsr + PLIC + IMEM + DMEM; AXI-Lite subsystem (axi\_interface, axi\_interconnect, 3×axi\_sfr); AHB-Lite subsystem (Request FIFO, Response FIFO, ahb\_interface, ahb\_interconnect, 3×ahb\_sfr); các luồng dữ liệu và điều khiển giữa các khối]}\vspace{5cm}}}
\caption{Sơ đồ khối tổng thể hệ thống SoC RISC-V}
\label{fig:soc-top-block}
\end{figure}

\subsection{Bản Đồ Bộ Nhớ}

CPU sử dụng không gian địa chỉ 32 bit (4\,GB). Bảng~\ref{tab:memory-map} mô tả các vùng địa chỉ được decode trong hệ thống.

\begin{table}[htbp]
\caption{Bản đồ bộ nhớ hệ thống SoC}
\label{tab:memory-map}
\begin{center}
\begin{tabular}{|l|l|l|l|l|}
\hline
\textbf{Vùng} & \textbf{Địa chỉ bắt đầu} & \textbf{Địa chỉ kết thúc} & \textbf{Kích thước} & \textbf{Bus/Latency} \\ \hline
IMEM & \texttt{0x0000\_0000} & \texttt{0x0000\_FFFF} & 64\,KB & Direct / 1 cycle \\ \hline
DMEM & \texttt{0x0001\_0000} & \texttt{0x0001\_FFFF} & 64\,KB & Direct / 1 cycle \\ \hline
PLIC & \texttt{0x0C00\_0000} & \texttt{0x0CFF\_FFFF} & 16\,MB & Direct / 1 cycle \\ \hline
AXI-Lite & \texttt{0x2000\_0000} & \texttt{0x2FFF\_FFFF} & 256\,MB & AXI / 4--6 cycle \\ \hline
AHB-Lite & \texttt{0x3000\_0000} & \texttt{0x3FFF\_FFFF} & 256\,MB & AHB+CDC / 6--12 cycle \\ \hline
\end{tabular}
\end{center}
\end{table}

Trong phạm vi vùng AXI và AHB, các slave được phân biệt bằng các bit địa chỉ \texttt{addr[27:12]} (slave ID) và \texttt{addr[11:0]} (offset trong slave). Việc này cho phép mở rộng tối đa 65536 slave mỗi bus và 4\,KB thanh ghi mỗi slave.

\subsection{Miền Xung Nhịp}

Hệ thống có hai miền xung nhịp:
\begin{itemize}
    \item \textbf{clk\_cpu (1\,GHz):} Toàn bộ CPU pipeline, IMEM, DMEM, PLIC, AXI interface, và Zicsr.
    \item \textbf{clk\_ahb (500\,MHz):} AHB interface FSM, ahb\_interconnect, và các AHB SFR slave.
\end{itemize}

Reset được thiết kế theo cơ chế \textbf{async assert, sync deassert}: khi mất reset, tín hiệu lan truyền ngay lập tức đến tất cả flip-flop; khi cấp lại reset, module \texttt{reset\_sync} cần 2 chu kỳ \texttt{clk\_ahb} để phục hồi an toàn tín hiệu \texttt{rst\_ahb\_n} cho domain 500\,MHz.

\section{Pipeline CPU 7 Tầng}

Pipeline CPU được tổ chức thành 7 tầng như minh họa trong Hình~\ref{fig:pipeline-block}. Giữa mỗi cặp tầng liên tiếp là một thanh ghi pipeline lưu trữ tất cả tín hiệu dữ liệu và điều khiển cần cho tầng kế tiếp.

\begin{figure}[htp]
\centering
\fbox{\parbox{0.85\textwidth}{\centering\textit{[Chèn hình: Sơ đồ pipeline chi tiết 7 tầng -- IF1 (PC), IF2 (IMEM read), ID (decode + regfile), EX (ALU + branch\_comp + addr\_adder), MEM1 (addr decode + bus dispatch), MEM2 (data mux + sign extension), WB (regfile write); các thanh ghi pipeline giữa các tầng; mũi tên forwarding từ MEM1/MEM2/WB về EX; mũi tên flush/stall từ hazard\_unit]}\vspace{4.5cm}}}
\caption{Sơ đồ chi tiết pipeline CPU 7 tầng}
\label{fig:pipeline-block}
\end{figure}

\subsection{Tầng IF1 -- Instruction Fetch Stage 1}

Tầng IF1 quản lý Program Counter (PC). Mỗi chu kỳ không bị stall, PC tăng thêm 4 (word-aligned fetch). Khi xảy ra branch taken hoặc trap/MRET, PC nhận địa chỉ mới từ \texttt{addr\_adder} (EX stage) hoặc từ \texttt{zicsr\_pc}. Ưu tiên: Zicsr > EX. PC được gửi trực tiếp đến IMEM trong cùng chu kỳ.

\subsection{Tầng IF2 -- Instruction Fetch Stage 2}

IMEM có latency 1 chu kỳ, nên kết quả đọc instruction xuất hiện ở tầng IF2. Tầng IF2 hội tụ (merge) giữa instruction từ IMEM và PC từ thanh ghi IF1/IF2, tạo ra cặp (PC, instr) chuyển sang ID. Module \texttt{imem} hỗ trợ tín hiệu \texttt{stall} (giữ output) và \texttt{flush} (xuất NOP -- \texttt{addi x0,x0,0}) để đồng bộ với pipeline.

\subsection{Tầng ID -- Instruction Decode}

Tầng ID bao gồm hai module: \texttt{id\_decoder} giải mã lệnh 32 bit thành tất cả tín hiệu điều khiển cần thiết, và \texttt{register\_file} đọc hai thanh ghi nguồn rs1/rs2. Module \texttt{id\_decoder} xác định loại lệnh, địa chỉ thanh ghi, immediate value, và các tín hiệu điều khiển cho các tầng sau.

Một điểm thiết kế quan trọng là \textbf{WBR bypass} trong \texttt{register\_file}: đọc thanh ghi kết hợp kiểm tra xem WB stage có đang ghi vào cùng địa chỉ không. Nếu có (và rd $\neq$ x0), dữ liệu từ WB được trả về trực tiếp thay vì đọc từ mảng thanh ghi. Điều này xử lý gap-4 RAW hazard mà không cần stall.

\subsection{Tầng EX -- Execute}

Tầng EX được đóng gói trong module \texttt{ex\_stage}, bao gồm:
\begin{itemize}
    \item \textbf{ALU:} Thực hiện 11 phép tính (ADD, SUB, AND, OR, XOR, SLL, SRL, SRA, SLT, SLTU, PASSB cho LUI).
    \item \textbf{Branch Comparator:} Đánh giá điều kiện branch (BEQ, BNE, BLT, BGE, BLTU, BGEU).
    \item \textbf{Address Adder:} Tính địa chỉ đích branch/jump (PC-relative hoặc register-relative cho JALR); JALR mask bit 0 theo đặc tả ISA.
    \item \textbf{Forwarding Unit và MUX:} Chọn giữa dữ liệu từ register file và forwarding từ MEM1/MEM2/WB.
\end{itemize}

\texttt{ex\_stage} xuất \texttt{branch\_taken} và \texttt{addr\_out} về hazard unit để quyết định flush pipeline khi branch taken hoặc jump.

\subsection{Tầng MEM1 -- Memory Access Stage 1}

Tầng MEM1 là điểm phân nhánh của địa chỉ: module \texttt{mem1\_stage} decode địa chỉ từ ALU result và phát tín hiệu điều khiển đến đúng subsystem:

\begin{itemize}
    \item \texttt{addr[31:16] == 16'h0001}: Truy cập DMEM -- không stall.
    \item \texttt{addr[31:24] == 8'h0C}: Truy cập PLIC -- không stall, latency 1 chu kỳ.
    \item \texttt{addr[31:28] == 4'h2}: Truy cập AXI-Lite -- phát \texttt{bus\_stall\_req = 1}.
    \item \texttt{addr[31:28] == 4'h3}: Truy cập AHB-Lite -- phát \texttt{bus\_stall\_req = 1}.
    \item Không khớp: Phát \texttt{load\_fault}/\texttt{store\_fault} để Zicsr sinh ngoại lệ.
\end{itemize}

Khi \texttt{bus\_stall\_req = 1}, toàn bộ pipeline từ IF1 đến MEM1 bị đóng băng (stall) cho đến khi bus transaction hoàn thành.

\subsection{Tầng MEM2 -- Memory Access Stage 2}

Tầng MEM2 thu kết quả từ DMEM (1 chu kỳ latency), AXI interface, hoặc AHB (qua Response FIFO). Tín hiệu \texttt{mem\_src[1:0]} chọn nguồn dữ liệu: DMEM, AXI, hoặc AHB. Tầng này cũng xử lý sign/zero extension cho các lệnh load byte và halfword (LB, LBU, LH, LHU) dựa trên \texttt{mem\_size} và \texttt{mem\_ext}.

\subsection{Tầng WB -- Write Back}

Tầng WB chọn dữ liệu kết quả cuối cùng và ghi vào register file. Tín hiệu \texttt{wb\_sel[1:0]} chọn từ bốn nguồn: kết quả ALU (00), dữ liệu load từ MEM2 (01), PC+4 cho JAL/JALR link (10), và dữ liệu CSR đọc (11).

\section{Xử Lý Hazard và Forwarding}

\subsection{Forwarding Unit}

Module \texttt{forwarding\_unit} phát hiện và giải quyết RAW hazard theo ưu tiên từ cao đến thấp: MEM1 (gap-1) $>$ MEM2 (gap-2) $>$ WB (gap-3). Kết quả là tín hiệu \texttt{fwd\_sel\_a[1:0]} và \texttt{fwd\_sel\_b[1:0]} chọn nguồn dữ liệu cho hai đầu vào của ALU. Nhờ forwarding, các lệnh RAW ở gap-1, gap-2, và gap-3 không cần stall.

\subsection{Hazard Unit}

Module \texttt{hazard\_unit} tổng hợp tất cả các tín hiệu stall và flush cho pipeline. Bảng~\ref{tab:hazard} tóm tắt các loại hazard và cách xử lý.

\begin{table}[htbp]
\caption{Tổng hợp các loại hazard và cơ chế xử lý trong pipeline 7 tầng}
\label{tab:hazard}
\begin{center}
\begin{tabular}{|l|l|c|}
\hline
\textbf{Loại Hazard} & \textbf{Cơ chế xử lý} & \textbf{Số chu kỳ stall} \\ \hline
Gap-1, 2, 3 RAW & Forwarding từ MEM1/MEM2/WB & 0 \\ \hline
Gap-4 RAW & WBR bypass trong register\_file & 0 \\ \hline
Load-use & Stall IF1..ID, bubble vào EX & 1 \\ \hline
CSR-use (EX) & Stall IF1..ID & 3 \\ \hline
CSR-use (MEM1) & Stall IF1..ID & 2 \\ \hline
CSR-use (MEM2) & Stall IF1..ID & 1 \\ \hline
Branch/Jump & Flush IF1/IF2 và IF2/ID & 2 (flush, không stall) \\ \hline
Bus stall (AXI/AHB) & Stall toàn pipeline IF1..MEM1 & N (đến khi xong) \\ \hline
Trap/MRET (Zicsr) & Flush toàn pipeline, load PC mới & -- \\ \hline
\end{tabular}
\end{center}
\end{table}

Một vấn đề thiết kế tinh tế là \textbf{bus stall kết hợp load-use hazard}: khi có lệnh SW ở MEM1 (đang bus stall), LW ở EX (có load-use với lệnh tiếp theo), và lệnh đọc ở ID -- ba điều kiện này có thể xảy ra đồng thời. Giải pháp đúng là khi \texttt{bus\_stall\_req = 1}, không được phát \texttt{flush\_id\_ex} cho load-use (vì toàn pipeline đã bị đóng băng, bubble sẽ được tạo tự nhiên sau khi bus giải phóng).

\section{Hệ Thống Bộ Nhớ}

IMEM và DMEM đều là bộ nhớ synchronous 1 chu kỳ latency, mỗi loại 64\,KB. Hệ thống sử dụng kiến trúc Harvard -- IMEM chỉ được đọc qua PC (instruction fetch path) và DMEM được truy cập qua ALU result (data path). Hai bộ nhớ này là hai mảng vật lý độc lập, nên không có structural hazard.

IMEM được tham số hóa với \texttt{SIZE\_KB} (mặc định 64\,KB) để cho phép compliance testbench override lên 512\,KB khi chạy bộ kiểm thử RV32I (một số bài kiểm thử branch có code size lên đến 293\,KB).

\section{Giao Diện Bus AXI-Lite}

\subsection{AXI Interface FSM}

Module \texttt{axi\_interface} là FSM (Finite State Machine) chuyển đổi từ giao thức nội bộ CPU (req\_valid/resp\_valid) sang AXI4-Lite 5 kênh. Đối với giao dịch \textbf{ghi}, FSM trải qua ba trạng thái: AW\_W\_PHASE (phát AWVALID và WVALID đồng thời), B\_PHASE (nhận BRESP), và trở về IDLE khi \texttt{axi\_resp\_valid = 1}. Đối với giao dịch \textbf{đọc}: AR\_PHASE (phát ARVALID), R\_PHASE (nhận RDATA), và trở về IDLE. Tín hiệu \texttt{axi\_resp\_valid = 1} báo hiệu \texttt{mem1\_stage} giải phóng \texttt{bus\_stall\_req}.

Hình~\ref{fig:axi-fsm} minh họa sơ đồ trạng thái của AXI Interface FSM.

\begin{figure}[htp]
\centering
\fbox{\parbox{0.85\textwidth}{\centering\textit{[Chèn hình: Sơ đồ FSM AXI Interface -- Write path: IDLE $\rightarrow$ AW\_W\_PHASE $\rightarrow$ B\_PHASE $\rightarrow$ IDLE; Read path: IDLE $\rightarrow$ AR\_PHASE $\rightarrow$ R\_PHASE $\rightarrow$ IDLE; điều kiện chuyển trạng thái trên từng mũi tên]}\vspace{3.5cm}}}
\caption{Sơ đồ trạng thái AXI Interface FSM}
\label{fig:axi-fsm}
\end{figure}

\subsection{AXI Interconnect và SFR Slave}

Module \texttt{axi\_interconnect} decode địa chỉ \texttt{addr[27:12]} để chọn một trong ba slave (S0, S1, S2). Mỗi slave là một \texttt{axi\_sfr} -- triển khai SFR Standard Register Map 9 thanh ghi (xem Mục~\ref{sec:sfr-map}). Trong trường hợp địa chỉ không decode được, interconnect trả về DECERR (BRESP/RRESP = 2'b11), kích hoạt bus error exception.

\section{Giao Diện Bus AHB-Lite và CDC}

\subsection{Asynchronous FIFO (CDC)}

Ranh giới giữa miền CPU (1\,GHz) và miền AHB (500\,MHz) được xử lý bởi hai FIFO bất đồng bộ:

\begin{itemize}
    \item \textbf{Request FIFO (1\,GHz $\rightarrow$ 500\,MHz):} Rộng 67 bit = Address(32) + WriteData(32) + Control(3: write/size). Được ghi bởi \texttt{mem1\_stage} và đọc bởi \texttt{ahb\_interface}.
    \item \textbf{Response FIFO (500\,MHz $\rightarrow$ 1\,GHz):} Rộng 33 bit = HRDATA(32) + HRESP(1). Được ghi bởi \texttt{ahb\_interface} và đọc bởi \texttt{mem1\_stage}.
\end{itemize}

Cả hai FIFO có độ sâu 2 (lũy thừa của 2, yêu cầu tối thiểu cho Gray code). Vì CPU bị stall cứng trong suốt bus transaction, không có overflow -- chỉ cần flag \texttt{empty} để kiểm tra trạng thái. Không dùng flag \texttt{full}.

\subsection{AHB Interface FSM}

Module \texttt{ahb\_interface} chạy ở 500\,MHz, đọc request từ Request FIFO và thực thi giao dịch AHB-Lite. FSM có ba trạng thái: IDLE (chờ request), ADDR (phát HADDR và HTRANS=NONSEQ), và DATA (phát HWDATA hoặc đọc HRDATA khi HREADY=1). Khi giao dịch hoàn thành, kết quả được ghi vào Response FIFO.

\subsection{Đồng Bộ IRQ AHB}

Tín hiệu ngắt từ các AHB slave (500\,MHz) cần được đồng bộ hóa trước khi vào PLIC (1\,GHz). Module \texttt{irq\_sync2ff} (2-FF synchronizer, 1\,GHz) thực hiện việc này -- có một instance cho mỗi AHB slave (tổng cộng ba instance trong \texttt{soc\_top}). IRQ từ AXI slave (1\,GHz) được kết nối thẳng vào PLIC.

\section{Bộ Điều Khiển Ngắt Ưu Tiên (PLIC)}

Module \texttt{plic} triển khai PLIC theo đặc tả SiFive với 6 nguồn ngắt (3 từ AXI slave + 3 từ AHB slave sau khi đồng bộ). Các tính năng chính:

\begin{itemize}
    \item \textbf{Priority per source:} Mỗi nguồn có thanh ghi ưu tiên 3 bit (0 = vô hiệu hóa, 1-7 = mức ưu tiên).
    \item \textbf{Threshold:} CPU có thể ghi ngưỡng ưu tiên; chỉ ngắt có ưu tiên cao hơn ngưỡng mới được chuyển lên.
    \item \textbf{Claim/Complete:} CPU đọc register CLAIM để nhận ID của ngắt đang chờ (ngắt có ưu tiên cao nhất, cao hơn ngưỡng); ghi vào COMPLETE sau khi xử lý xong.
    \item \textbf{Latency 1 chu kỳ:} PLIC được ánh xạ vào địa chỉ \texttt{0x0C000000}, phản hồi trong 1 chu kỳ, không gây stall.
\end{itemize}

PLIC xuất tín hiệu \texttt{meip\_in} đến module Zicsr để kích hoạt Machine External Interrupt.

\section{Khối CSR và Xử Lý Ngắt/Ngoại Lệ (Zicsr)}

Module \texttt{zicsr} quản lý sáu thanh ghi CSR và xử lý toàn bộ trap (ngắt + ngoại lệ) tại M-mode. Zicsr nhận tín hiệu từ tầng WB của pipeline (vì ngắt và ngoại lệ được xử lý khi lệnh gây ra nó đến WB -- đây là điều kiện \textbf{precise exception}).

\subsection{Các Nguồn Trap}

\begin{itemize}
    \item \textbf{Exception:} ECALL, EBREAK, Illegal Instruction, Load/Store Access Fault -- khi lệnh đến WB.
    \item \textbf{Interrupt:} Machine External Interrupt (từ PLIC), Machine Software Interrupt (từ phần mềm) -- được kiểm tra mỗi chu kỳ khi \texttt{mstatus.MIE = 1}.
    \item \textbf{Bus Error:} BRESP/RRESP SLVERR (AXI) hoặc HRESP ERROR (AHB) -- được phát hiện và chuyển thành Load/Store Fault.
\end{itemize}

\subsection{Precise Exception}

Khi có bus transaction đang diễn ra (\texttt{bus\_stall\_req = 1}), Zicsr \textbf{không} được flush pipeline. Điều này đảm bảo bus transaction không bị hủy ngang chừng -- \textbf{precise exception} đòi hỏi phải chờ giao dịch bus hoàn thành trước khi xử lý trap.

\subsection{Vectored Interrupt Mode}

Với \texttt{mtvec[1:0] = 01} (vectored mode), địa chỉ vector trap được tính là:
\begin{equation}
    \text{PC}_{\text{trap}} = (\texttt{mtvec} \, \& \, \sim3) + 4 \times \texttt{cause\_code}
\end{equation}
Mỗi nguyên nhân trap có vector riêng biệt, giúp giảm độ trễ xử lý ngắt bằng cách tránh switch-case trong phần mềm.

\section{SFR Standard Register Map}
\label{sec:sfr-map}

Mọi ngoại vi kết nối vào hệ thống phải triển khai SFR Standard Register Map (lấy cảm hứng từ OpenTitan), đảm bảo tính nhất quán và khả năng plug-and-play. Bảng~\ref{tab:sfr-map} mô tả cấu trúc thanh ghi chuẩn.

\begin{table}[htbp]
\caption{SFR Standard Register Map -- bố cục thanh ghi chuẩn cho mọi ngoại vi}
\label{tab:sfr-map}
\begin{center}
\begin{tabular}{|c|l|l|l|}
\hline
\textbf{Offset} & \textbf{Tên} & \textbf{Access} & \textbf{Mô tả} \\ \hline
\texttt{0x00} & CTRL & RW & \texttt{bit[0]} = enable; bits[31:1] = peripheral-specific \\ \hline
\texttt{0x04} & STATUS & RO & Trạng thái read-only do peripheral drive \\ \hline
\texttt{0x08} & INTR\_ENABLE & RW & Mask từng nguồn IRQ \\ \hline
\texttt{0x0C} & INTR\_STATE & RW1C & Pending flags -- ghi 1 để clear \\ \hline
\texttt{0x10} & INTR\_TEST & WO & Ghi 1 để force-set INTR\_STATE (debug) \\ \hline
\texttt{0x14} & DATA0 & RW & General-purpose (peripheral-specific) \\ \hline
\texttt{0x18} & DATA1 & RW & General-purpose \\ \hline
\texttt{0x1C} & DATA2 & RW & General-purpose \\ \hline
\texttt{0xFC} & PERIPH\_ID & RO & Hardcoded peripheral identifier \\ \hline
\end{tabular}
\end{center}
\end{table}

Quy tắc IRQ: \texttt{irq = |(INTR\_STATE \& INTR\_ENABLE)} -- ngắt chỉ được phát khi có event đang pending VÀ được enable. Thanh ghi INTR\_TEST cho phép phần mềm kiểm tra đường ngắt mà không cần kích hoạt sự kiện thực sự.

Trong khóa luận này, hai loại SFR được triển khai: \texttt{axi\_sfr} cho bus AXI-Lite và \texttt{ahb\_sfr} cho bus AHB-Lite. Ngoài ra, \texttt{gpio\_sfr} là một ví dụ ngoại vi hoàn chỉnh: \texttt{DATA0} điều khiển GPIO output, \texttt{STATUS} phản ánh GPIO input, và edge-detect circuit tạo IRQ khi có cạnh tín hiệu.

\section{Kiến Trúc Chip Boundary}

Module \texttt{soc\_top} được thiết kế như một \textbf{chip boundary}: tất cả các slave AXI và AHB được expose ra ngoài thông qua port declarations, cho phép kết nối với ngoại vi bên ngoài mà không cần sửa đổi RTL bên trong chip. \texttt{soc\_top} xuất:

\begin{itemize}
    \item \texttt{rst\_cpu\_n\_o}, \texttt{rst\_ahb\_n\_o}: Reset đã đồng bộ, để ngoại vi dùng reset flip-flop của mình.
    \item \texttt{axi\_S\{0,1,2\}\_*}: Đầy đủ cổng AXI-Lite slave cho ba slave.
    \item \texttt{ahb\_S\{0,1,2\}\_*}: Cổng AHB-Lite cho ba slave.
\end{itemize}

Triết lý này đảm bảo bất kỳ ngoại vi nào tuân thủ SFR Standard đều có thể được gắn vào mà không cần thay đổi RTL trong chip.


\chapter{KẾT QUẢ THỰC NGHIỆM}
\label{Chapter5}

\section{Môi Trường Kiểm Thử}

\subsection{Công Cụ và Thiết Lập}

Toàn bộ quá trình kiểm thử được thực hiện trên hệ thống Linux với các công cụ sau:

\begin{itemize}
    \item \textbf{Icarus Verilog 12} (\texttt{iverilog -g2012 -Wall}): Biên dịch và mô phỏng RTL SystemVerilog.
    \item \textbf{riscv64-unknown-elf-gcc 13.2.0}: Biên dịch chương trình thử nghiệm RISC-V.
    \item \textbf{GTKWave}: Xem và phân tích dạng sóng VCD.
    \item \textbf{GNU Make}: Tự động hóa toàn bộ quy trình build và test.
    \item \textbf{riscv-arch-test (old-framework-2.x)}: Bộ kiểm thử tuân thủ RV32I chính thức.
\end{itemize}

Mọi testbench đều sử dụng \textbf{tự kiểm tra tự động} (self-checking): tín hiệu \texttt{PASS}/\texttt{FAIL} được in ra sau mỗi test case mà không cần so sánh thủ công.

\subsection{Chiến Lược Kiểm Thử}

Chiến lược kiểm thử được tổ chức thành các pha tăng dần về độ phức tạp:

\begin{figure}[htp]
\centering
\fbox{\parbox{0.85\textwidth}{\centering\textit{[Chèn hình: Tháp kiểm thử hình kim tự tháp -- từ dưới lên: Unit Test (ALU, branch\_comp, register\_file, id\_decoder, forwarding\_unit, hazard\_unit, async\_fifo, plic, ex\_stage, irq\_sync2ff, gpio\_sfr, zicsr), Integration Test (axi\_interface, ahb\_interface, axi\_full, ahb\_full, bus\_error, ahb\_error), System Test (tb\_pipeline\_cpu, tb\_soc\_top 20 programs), Compliance Test (riscv-arch-test 38 tests)]}\vspace{3.5cm}}}
\caption{Chiến lược kiểm thử phân cấp}
\label{fig:test-pyramid}
\end{figure}

\section{Kiểm Thử Đơn Vị (Phase 1 và 2)}

\subsection{Phase 1 -- Kiểm Thử Module Tổ Hợp}

Phase 1 kiểm thử bốn module logic tổ hợp và đồng bộ quan trọng nhất của pipeline. Tổng cộng 192 test case, tất cả PASS.

\subsubsection{ALU -- 38 Test Case}

Module ALU (\texttt{alu.sv}) thực hiện 11 phép tính với các giá trị biên đặc trưng. Bảng~\ref{tab:alu-tests} tóm tắt kết quả.

\begin{table}[htbp]
\caption{Kết quả kiểm thử ALU -- 38 test case}
\label{tab:alu-tests}
\begin{center}
\begin{tabular}{|l|c|l|}
\hline
\textbf{Nhóm phép tính} & \textbf{Số test} & \textbf{Corner case tiêu biểu} \\ \hline
ADD & 4 & Overflow wrap: \texttt{0xFFFF\_FFFF + 1 = 0} \\ \hline
SUB & 3 & Underflow: \texttt{0 - 1 = 0xFFFF\_FFFF} \\ \hline
SLL/SRL/SRA & 10 & Shift 0, shift 31, shift 32 (modulo 32) \\ \hline
SLT/SLTU & 5 & So sánh số âm và số dương, zero \\ \hline
AND/OR/XOR & 8 & All zeros, all ones, alternating bits \\ \hline
PASSB (LUI) & 4 & Chuyển nguyên immediate, kiểm tra upper bits \\ \hline
\multicolumn{2}{|c|}{\textbf{Tổng}} & \textbf{38/38 PASS} \\ \hline
\end{tabular}
\end{center}
\end{table}

\subsubsection{Branch Comparator -- 30 Test Case}

Module \texttt{branch\_comp} kiểm thử sáu điều kiện branch: BEQ, BNE, BLT, BGE, BLTU, BGEU. Đặc biệt chú ý đến trường hợp biên: số âm lớn nhất, số dương lớn nhất, và so sánh có/không có dấu (BLT vs BLTU cho giá trị 0x80000000).

\subsubsection{Register File -- 36 Test Case}

Kiểm thử \texttt{register\_file} tập trung vào hai tính năng đặc biệt: (1) x0 hardwired zero (không thể ghi), (2) WBR bypass (gap-4 RAW hazard) -- đọc và ghi cùng địa chỉ trong cùng chu kỳ phải trả về dữ liệu mới.

\subsubsection{Instruction Decoder -- 88 Test Case}

Module \texttt{id\_decoder} được kiểm thử với tất cả loại lệnh RV32I và Zicsr. Kiểm tra đặc biệt: immediate sign-extension cho từng định dạng lệnh, xác định đúng \texttt{alu\_op} cho từng lệnh.

\subsection{Phase 2 -- Kiểm Thử Module Đồng Bộ}

Phase 2 bổ sung ba module có trạng thái (stateful): \texttt{forwarding\_unit}, \texttt{hazard\_unit}, và \texttt{async\_fifo}. Tổng cộng 114 test case, tất cả PASS.

\subsubsection{Forwarding Unit -- 31 Test Case}

Kiểm thử tất cả tổ hợp forwarding cho rs1 và rs2: không forward (00), từ MEM1 (01), từ MEM2 (10), từ WB (11). Đặc biệt kiểm tra trường hợp rd = x0 (không được forward), và conflict giữa forwarding từ MEM1 và MEM2 (MEM1 phải thắng).

\subsubsection{Hazard Unit -- 73 Test Case}

Module \texttt{hazard\_unit} là module phức tạp nhất trong pipeline. Các test case bao gồm:
\begin{itemize}
    \item Load-use hazard đơn giản (1 stall cycle).
    \item CSR-use hazard ở gap 1, 2, 3 (3/2/1 stall cycle tương ứng).
    \item Branch/Jump flush (2 slot flush).
    \item Bus stall kết hợp load-use -- tình huống đặc biệt quan trọng.
    \item Zicsr flush toàn pipeline.
\end{itemize}

\subsubsection{Async FIFO -- 10 Test Case}

Module \texttt{async\_fifo} (depth=2) được kiểm thử với các trường hợp: ghi và đọc bình thường, kiểm tra flag empty chính xác, ghi khi FIFO trống, đọc khi FIFO có 1 phần tử, và kiểm tra Gray code pointer ở biên (wrap-around).

\section{Kiểm Thử Tích Hợp Pipeline (Phase 3)}

Phase 3 là lần đầu tiên toàn bộ pipeline CPU được chạy với chương trình thực sự thông qua \texttt{soc\_top}. Chín chương trình thử nghiệm kiểm tra từng nhóm lệnh, từng loại hazard và các tình huống đặc biệt. Kết quả: \textbf{9/9 PASS}.

\begin{table}[htbp]
\caption{Chương trình kiểm thử pipeline (Phase 3) -- 9/9 PASS}
\label{tab:phase3-programs}
\begin{center}
\begin{tabular}{|l|l|}
\hline
\textbf{Chương trình} & \textbf{Nội dung kiểm thử} \\ \hline
prog\_arithmetic & Tất cả lệnh số học: ADD/SUB/AND/OR/XOR/SLT/LUI/AUIPC \\ \hline
prog\_shift & SLL/SRL/SRA với shift amount biên (0, 1, 31) \\ \hline
prog\_forwarding & RAW hazard gap-1, 2, 3, 4 với nhiều chuỗi lệnh phụ thuộc \\ \hline
prog\_load\_store & LW/LH/LB/LHU/LBU/SW/SH/SB -- địa chỉ và dữ liệu biên \\ \hline
prog\_branch\_jump & BEQ/BNE/BLT/BGE/BLTU/BGEU taken và not-taken; JAL/JALR \\ \hline
prog\_load\_use & Load-use hazard: LW theo sau ngay bởi lệnh dùng kết quả \\ \hline
prog\_ecall & ECALL, EBREAK, Illegal Instruction → trap handler → MRET \\ \hline
prog\_csr & CSRRW/CSRRS/CSRRC/CSRxI -- đọc/ghi/modify mstatus/mie/mtvec \\ \hline
prog\_interrupt & Software interrupt (MSI) -- kích hoạt, xử lý handler, clear \\ \hline
\end{tabular}
\end{center}
\end{table}

Trong quá trình kiểm thử Phase 3, hai bug quan trọng trong RTL được phát hiện và sửa:

\begin{enumerate}
    \item \textbf{Bug B3 -- Gap-4 RAW hazard:} \texttt{register\_file} không có WBR bypass cho trường hợp WB ghi và ID đọc cùng địa chỉ trong cùng chu kỳ. Sửa bằng cách thêm combinational bypass.
    \item \textbf{Bug B6 -- Ghost instruction:} Sau \texttt{zicsr\_flush}, IMEM tiếp tục xuất instruction của địa chỉ cũ trong chu kỳ tiếp theo, gây illegal instruction spurious. Sửa bằng cách thêm tín hiệu flush vào IMEM để xuất NOP.
\end{enumerate}

\section{Kiểm Thử Giao Diện Bus (Phase 4)}

Phase 4 được chia thành bốn sub-phase, tăng dần từ kiểm thử interface đơn lẻ đến full path với interconnect và SFR slave. Tổng cộng 163 test case, tất cả PASS.

\subsection{Phase 4a -- AXI Interface (49 Test Case)}

\texttt{tb\_axi\_interface} sử dụng slave model AXI giả để kiểm thử module \texttt{axi\_interface}. Các test case bao gồm:

\begin{itemize}
    \item Giao dịch ghi: one-shot, wait state (slave AWREADY delay), BRESP SLVERR (bus error).
    \item Giao dịch đọc: one-shot, ARREADY delay, RVALID delay, RRESP SLVERR.
    \item Kiểm tra bus stall: pipeline giữ nguyên trong suốt giao dịch.
    \item Kiểm tra \texttt{axi\_resp\_valid} timing: bus stall giải phóng đúng thời điểm.
\end{itemize}

\subsection{Phase 4b -- AHB Interface và CDC (29 Test Case)}

\texttt{tb\_ahb\_interface} kiểm thử \texttt{ahb\_interface} kết hợp với hai FIFO bất đồng bộ. Đây là kiểm thử CDC đầu tiên -- các giao dịch được phát từ domain 1\,GHz và nhận ở domain 500\,MHz, sau đó kết quả đi ngược về 1\,GHz. Test case đặc biệt: HRESP ERROR (AHB bus error) → phát hiện và báo lên CPU.

Hình~\ref{fig:cdc-waveform} minh họa waveform của một giao dịch hoàn chỉnh qua CDC.

\begin{figure}[htp]
\centering
\fbox{\parbox{0.85\textwidth}{\centering\textit{[Chèn hình: Waveform CDC -- ghi request vào Request FIFO (clk\_cpu), đọc request từ FIFO (clk\_ahb), thực thi AHB transaction (HADDR, HWDATA, HREADY), ghi response vào Response FIFO (clk\_ahb), đọc response từ FIFO (clk\_cpu), giải phóng bus\_stall\_req]}\vspace{4cm}}}
\caption{Dạng sóng giao dịch CDC -- Request FIFO và Response FIFO}
\label{fig:cdc-waveform}
\end{figure}

\subsection{Phase 4c -- AXI Full Path (47 Test Case)}

\texttt{tb\_axi\_full} kiểm thử chuỗi đầy đủ: \texttt{axi\_interface} + \texttt{axi\_interconnect} + 3~\texttt{axi\_sfr}. Các test case bao gồm:
\begin{itemize}
    \item Ghi/đọc đến từng slave S0, S1, S2.
    \item Kiểm tra decode địa chỉ đúng: ghi vào S0 không ảnh hưởng S1, S2.
    \item Kiểm tra toàn bộ 9 thanh ghi SFR Standard per slave.
    \item Kiểm tra INTR\_STATE/INTR\_ENABLE và logic IRQ.
    \item Kiểm tra INTR\_TEST: force-set INTR\_STATE để kiểm tra đường IRQ.
\end{itemize}

\subsection{Phase 4d -- AHB Full Path (38 Test Case)}

\texttt{tb\_ahb\_full} là tương đương của Phase 4c cho bus AHB: \texttt{ahb\_interface} + CDC FIFO + \texttt{ahb\_interconnect} + 3~\texttt{ahb\_sfr}. Kiểm tra wait state (HREADYOUT = 0 từ slave), HRESP ERROR propagation, và SFR Standard Register Map qua AHB protocol.

\section{Kiểm Thử Hệ Thống (Phase 5 và 6)}

\subsection{Phase 5 -- Tích Hợp SoC + AXI/AHB (4 Chương Trình)}

Phase 5 là lần đầu tiên CPU thực sự truy cập ngoại vi qua bus thông qua \texttt{soc\_top}. Bốn chương trình kiểm thử:

\begin{enumerate}
    \item \textbf{prog\_axi\_sfr}: CPU ghi/đọc các thanh ghi AXI SFR; kiểm tra kết quả đúng.
    \item \textbf{prog\_ahb\_sfr}: CPU ghi/đọc các thanh ghi AHB SFR qua CDC; kiểm tra kết quả đúng.
    \item \textbf{prog\_axi\_irq}: CPU cấu hình AXI SFR IRQ, kích hoạt ngắt qua INTR\_TEST, xử lý ISR, clear INTR\_STATE.
    \item \textbf{prog\_ahb\_irq}: Tương tự với AHB SFR -- IRQ từ 500\,MHz domain, qua 2-FF sync, đến PLIC, đến CPU.
\end{enumerate}

Kết quả: \textbf{4/4 PASS}. Trong quá trình này, bug B8 được phát hiện: bus stall kết hợp load-use hazard gây LW bị hủy oan. Sửa trong \texttt{hazard\_unit.sv}.

\subsection{Phase 6a -- System Batch Test (20 Chương Trình)}

\texttt{tb\_soc\_top} chạy tự động 20 chương trình qua \texttt{soc\_top} với sequence reset chuẩn, bao gồm tất cả 9 chương trình từ Phase 3 cộng thêm 11 chương trình mới kiểm tra các tình huống phức tạp hơn.

\begin{figure}[htp]
\centering
\fbox{\parbox{0.85\textwidth}{\centering\textit{[Chèn hình: Waveform pipeline hoạt động -- 5-10 chu kỳ liên tiếp thể hiện pipeline đầy đủ 7 tầng: IF1 (PC), IF2 (instr từ IMEM), ID (decode), EX (ALU result), MEM1/MEM2/WB; forwarding paths active; có 1 stall cycle do load-use; flush do branch taken]}\vspace{4cm}}}
\caption{Dạng sóng pipeline CPU trong quá trình mô phỏng hệ thống}
\label{fig:pipeline-waveform}
\end{figure}

Kết quả: \textbf{20/20 PASS}.

\subsection{Phase 6b -- Compliance Framework (3 Chương Trình)}

\texttt{tb\_compliance} sử dụng framework riêng tương tự riscv-arch-test để kiểm tra các trường hợp đặc biệt: shifts (tất cả shift operations với biên), compare (SLT/SLTU với signed/unsigned edge cases), và dmem\_endurance (ghi/đọc 256 địa chỉ DMEM liên tiếp). Kết quả: \textbf{3/3 TEST\_PASS}.

\section{Kiểm Thử PLIC và EX Stage (Phase 7 và 8)}

\subsection{Phase 7 -- PLIC Unit Test (31 Test Case) và Hệ Thống (3 Chương Trình)}

\texttt{tb\_plic} kiểm thử module \texttt{plic} với 31 test case bao gồm:
\begin{itemize}
    \item Không có ngắt -- PLIC xuất meip=0.
    \item Ngắt đơn nguồn với priority > 0.
    \item Nhiều ngắt đồng thời -- PLIC chọn nguồn có priority cao nhất.
    \item Threshold: ngắt có priority = threshold không được chuyển lên CPU.
    \item Claim/Complete cycle: CLAIM trả về ID đúng, COMPLETE clear pending.
    \item Edge case: priority = 0 (disabled), priority wrap-around.
\end{itemize}

Ba chương trình hệ thống (\texttt{prog\_plic\_basic}, \texttt{prog\_plic\_priority}, \texttt{prog\_plic\_threshold}) kiểm tra PLIC trong môi trường SoC hoàn chỉnh. Kết quả: \textbf{31/31 PASS + 3/3 PASS}.

\subsection{Phase 8 -- EX Stage Unit Test (23 Test Case) và CSR Hazard}

\texttt{tb\_ex\_stage} kiểm thử module \texttt{ex\_stage} như một đơn vị hoàn chỉnh (bao gồm forwarding\_unit, ALU, branch\_comp, addr\_adder và các MUX bên trong). Đặc biệt kiểm tra:
\begin{itemize}
    \item Forwarding override: khi có forwarding, ALU sử dụng dữ liệu mới nhất.
    \item Branch resolution: branch\_taken đúng với tất cả sáu điều kiện.
    \item JALR address masking: bit 0 của địa chỉ đích được mask.
\end{itemize}

Chương trình \texttt{prog\_csr\_hazard} chạy qua tất cả CSR-use gap (0--4): lệnh CSR ở EX/MEM1/MEM2/WB và lệnh tiếp theo đọc kết quả CSR với các khoảng cách khác nhau. Hazard unit phải stall đúng số chu kỳ. Kết quả: \textbf{23/23 PASS + 1/1 PASS}.

\section{Kiểm Thử Đơn Vị Bổ Sung}

Ba module phụ được bổ sung unit test riêng trong giai đoạn hoàn thiện:

\begin{table}[htbp]
\caption{Kết quả kiểm thử đơn vị bổ sung}
\label{tab:new-unit-tests}
\begin{center}
\begin{tabular}{|l|c|l|}
\hline
\textbf{Testbench} & \textbf{Kết quả} & \textbf{Nội dung chính} \\ \hline
tb\_irq\_sync2ff & 10/10 PASS & 2-FF sync: glitch rejection, latency 2 cycle \\ \hline
tb\_gpio\_sfr & 22/22 PASS & GPIO output (DATA0), input (STATUS), edge-detect IRQ \\ \hline
tb\_zicsr & 38/38 PASS & CSRRW/RS/RC, trap handler, MRET, bus error exception \\ \hline
\end{tabular}
\end{center}
\end{table}

\section{Kiểm Thử Bus Error}

Hai testbench tích hợp kiểm tra trường hợp bus trả lỗi:

\begin{itemize}
    \item \textbf{tb\_soc\_bus\_err:} AXI slave trả BRESP=SLVERR (store) → \texttt{mcause=7} (Store Fault); RRESP=SLVERR (load) → \texttt{mcause=5} (Load Fault). Kết quả: \textbf{2/2 PASS}.
    \item \textbf{tb\_soc\_ahb\_err:} AHB slave trả HRESP=ERROR → load/store fault tương ứng. Kết quả: \textbf{2/2 PASS}.
\end{itemize}

Hình~\ref{fig:bus-error-waveform} minh họa quá trình phát hiện và xử lý bus error.

\begin{figure}[htp]
\centering
\fbox{\parbox{0.85\textwidth}{\centering\textit{[Chèn hình: Waveform bus error -- AXI giao dịch ghi, slave trả BRESP=SLVERR (2'b10), axi\_interface báo axi\_resp\_err=1, mem2\_stage truyền store\_fault lên Zicsr, Zicsr flush pipeline và load mtvec, PC nhảy đến handler, mcause=7 được ghi]}\vspace{4cm}}}
\caption{Dạng sóng xử lý bus error (BRESP SLVERR $\rightarrow$ Store Fault)}
\label{fig:bus-error-waveform}
\end{figure}

\section{Kiểm Thử Tuân Thủ RV32I}

\subsection{Bộ Kiểm Thử riscv-arch-test}

Bộ kiểm thử tuân thủ chính thức \texttt{riscv-arch-test} (old-framework-2.x) của tổ chức RISC-V International \cite{riscv-arch-test} kiểm tra tất cả lệnh cơ bản của RV32I bằng phương pháp signature-based:

\begin{enumerate}
    \item Mỗi test program thực thi một chuỗi lệnh và ghi kết quả vào vùng nhớ ``signature'' trong DMEM.
    \item Sau khi mô phỏng, testbench dump signature ra file.
    \item Script so sánh signature với file tham chiếu được tạo từ ISS chuẩn (Spike).
    \item Nếu signature khớp hoàn toàn: PASS; ngược lại: FAIL.
\end{enumerate}

\subsection{Thiết Lập Đặc Biệt}

Một số điều chỉnh kỹ thuật cần thiết để chạy bộ kiểm thử này:

\textbf{1. IMEM 512\,KB cho branch tests:} Các test BEQ/BNE/BLT/BGE/BLTU/BGEU sử dụng test vector rất dày đặc, sinh code size lên đến 293\,KB -- vượt quá IMEM 64\,KB của phần cứng. Giải pháp: tham số hóa IMEM (\texttt{IMEM\_SIZE\_KB}) và override lên 512\,KB trong compliance testbench. Đây là kiến trúc Harvard nên IMEM và DMEM là hai mảng độc lập -- không có xung đột địa chỉ.

\textbf{2. DMEM preload cho load tests:} Các test LB/LBU/LH/LHU/LW đọc dữ liệu từ vùng \texttt{rvtest\_data} được khai báo trong section \texttt{.data}. Testbench cần đọc section này từ file ELF và nạp vào DMEM trước khi chạy mô phỏng.

\textbf{3. X-bit handling:} Icarus Verilog khởi tạo mảng DMEM với giá trị X (không xác định). Các từ trong signature chưa được ghi phải xuất ra \texttt{"00000000"} (Spike khởi tạo bộ nhớ bằng 0). Kiểm tra \texttt{\^{}word === 1'bx} để phân biệt từ chứa X.

\textbf{4. Patch jalr-01.S:} Một test case trong \texttt{jalr-01.S} dùng macro mở rộng thành lệnh \texttt{la x0, label} -- bị binutils $\geq$2.39 từ chối vì x0 là zero register. Patch thay bằng code tương đương dùng register hợp lệ, cho cùng kết quả signature.

\subsection{Kết Quả Tuân Thủ}

Bảng~\ref{tab:compliance-results} trình bày kết quả chi tiết của bộ kiểm thử RV32I.

\begin{table}[htbp]
\caption{Kết quả kiểm thử tuân thủ RV32I -- riscv-arch-test}
\label{tab:compliance-results}
\begin{center}
\begin{tabular}{|l|c|l|}
\hline
\textbf{Nhóm lệnh} & \textbf{Kết quả} & \textbf{Ghi chú} \\ \hline
addi-01 & PASS & \\ \hline
add-01 & PASS & \\ \hline
andi-01 & PASS & \\ \hline
and-01 & PASS & \\ \hline
auipc-01 & PASS & \\ \hline
beq-01 & PASS & 512\,KB IMEM (code 220\,KB) \\ \hline
bge-01 & PASS & 512\,KB IMEM (code 220\,KB) \\ \hline
bgeu-01 & PASS & 512\,KB IMEM (code 291\,KB) \\ \hline
blt-01 & PASS & 512\,KB IMEM (code 220\,KB) \\ \hline
bltu-01 & PASS & 512\,KB IMEM (code 293\,KB) \\ \hline
bne-01 & PASS & 512\,KB IMEM (code 220\,KB) \\ \hline
jal-01 & SKIP & Code $\approx$1.7\,MB, cần IMEM $>$2\,MB \\ \hline
jalr-01 & PASS & Patch test case inst\_7 (rd=x0 binutils bug) \\ \hline
lb-align-01 & PASS & DMEM preload từ .data section \\ \hline
lbu-align-01 & PASS & \\ \hline
lh-align-01 & PASS & \\ \hline
lhu-align-01 & PASS & \\ \hline
lui-01 & PASS & \\ \hline
lw-align-01 & PASS & \\ \hline
ori-01 & PASS & \\ \hline
or-01 & PASS & \\ \hline
sb-align-01 & PASS & \\ \hline
sh-align-01 & PASS & \\ \hline
slli-01 & PASS & \\ \hline
sll-01 & PASS & \\ \hline
slti-01 & PASS & \\ \hline
slt-01 & PASS & \\ \hline
sltiu-01 & PASS & \\ \hline
sltu-01 & PASS & \\ \hline
srai-01 & PASS & \\ \hline
sra-01 & PASS & \\ \hline
srli-01 & PASS & \\ \hline
srl-01 & PASS & \\ \hline
sub-01 & PASS & \\ \hline
sw-align-01 & PASS & \\ \hline
xori-01 & PASS & \\ \hline
xor-01 & PASS & \\ \hline
\hline
\multicolumn{2}{|c|}{\textbf{Tổng}} & \textbf{37/38 PASS; 1 SKIP (jal-01)} \\ \hline
\end{tabular}
\end{center}
\end{table}

Test jal-01 được phân loại là SKIP (không phải FAIL): bài kiểm thử này thiết kế để test JAL với offset $\pm$1\,MB, đòi hỏi code size khoảng 1.7\,MB -- vượt quá giới hạn thực tế của bất kỳ IMEM nhúng nào. Lệnh JAL đã được kiểm chứng hoạt động đúng thông qua 37 chương trình system test (mọi function call đều sử dụng JAL).

\section{Tổng Kết và Đánh Giá}

\subsection{Tổng Hợp Kết Quả}

Bảng~\ref{tab:test-summary} tổng hợp toàn bộ kết quả kiểm thử.

\begin{table}[htbp]
\caption{Tổng hợp kết quả kiểm thử toàn dự án}
\label{tab:test-summary}
\begin{center}
\begin{tabular}{|l|l|c|c|}
\hline
\textbf{Pha} & \textbf{Mô tả} & \textbf{Số lượng} & \textbf{Kết quả} \\ \hline
Phase 1 & Unit test: alu, branch\_comp, regfile, id\_decoder & 192 cases & 192/192 PASS \\ \hline
Phase 2 & Unit test: forwarding, hazard, async\_fifo & 114 cases & 114/114 PASS \\ \hline
Phase 3 & Integration: pipeline đầy đủ (9 programs) & 9 programs & 9/9 PASS \\ \hline
Phase 4a & Integration: AXI interface & 49 cases & 49/49 PASS \\ \hline
Phase 4b & Integration: AHB interface + CDC & 29 cases & 29/29 PASS \\ \hline
Phase 4c & Integration: AXI full path + 3 SFR & 47 cases & 47/47 PASS \\ \hline
Phase 4d & Integration: AHB full path + 3 SFR & 38 cases & 38/38 PASS \\ \hline
Phase 5 & System: SoC + AXI/AHB SFR + IRQ (4 programs) & 4 programs & 4/4 PASS \\ \hline
Phase 6a & System: batch 20 programs qua soc\_top & 20 programs & 20/20 PASS \\ \hline
Phase 6b & Compliance framework (3 programs) & 3 programs & 3/3 PASS \\ \hline
Phase 7 & Unit: PLIC (31) + system (3 programs) & 31+3 & 31/31+3/3 PASS \\ \hline
Phase 8 & Unit: EX stage (23) + CSR hazard (1) & 23+1 & 23/23+1/1 PASS \\ \hline
New unit & irq\_sync2ff (10), gpio\_sfr (22), zicsr (38) & 70 cases & 70/70 PASS \\ \hline
Bus error & AXI bus error + AHB bus error & 4 cases & 4/4 PASS \\ \hline
Compliance & riscv-arch-test RV32I & 38 tests & 37 PASS, 1 SKIP \\ \hline
\hline
\multicolumn{2}{|c|}{\textbf{Tổng}} & \textbf{$>$630 cases + 37 programs} & \textbf{Tất cả PASS} \\ \hline
\end{tabular}
\end{center}
\end{table}

\subsection{Bug Phát Hiện và Sửa Chữa}

Trong quá trình kiểm thử, tám bug RTL và testbench đã được phát hiện và sửa chữa. Bảng~\ref{tab:bugs} liệt kê các bug RTL quan trọng nhất.

\begin{table}[htbp]
\caption{Các bug RTL quan trọng phát hiện trong quá trình kiểm thử}
\label{tab:bugs}
\begin{center}
\begin{tabular}{|c|l|l|l|}
\hline
\textbf{\#} & \textbf{Bug} & \textbf{Module} & \textbf{Tác động} \\ \hline
B3 & Gap-4 RAW hazard thiếu WBR bypass & register\_file & Dữ liệu sai \\ \hline
B4 & IMEM synchronous stall mismatch & imem & Instruction bị mất \\ \hline
B5 & CSR-use hazard stall quá muộn & hazard\_unit & CSR read sai \\ \hline
B6 & Ghost instruction sau zicsr\_flush & imem & Vòng lặp trap vô hạn \\ \hline
B8 & bus\_stall + load\_use → LW bị hủy & hazard\_unit & Bus read không xảy ra \\ \hline
\end{tabular}
\end{center}
\end{table}

\subsection{Phân Tích Hiệu Năng Pipeline}

Trong điều kiện lý tưởng (không hazard), pipeline 7 tầng đạt throughput 1 lệnh/chu kỳ (IPC = 1). Các trường hợp làm giảm IPC trong thực tế:

\begin{itemize}
    \item \textbf{Branch/Jump:} 2 chu kỳ flush mỗi lần rẽ nhánh (không có branch prediction).
    \item \textbf{Load-use:} 1 chu kỳ stall mỗi lần LW theo sau ngay bởi lệnh dùng kết quả.
    \item \textbf{CSR-use:} 1--3 chu kỳ stall tùy khoảng cách.
    \item \textbf{Bus stall:} 4--6 chu kỳ cho AXI, 6--12 chu kỳ cho AHB (bao gồm CDC latency).
\end{itemize}

Nhờ forwarding unit hoàn chỉnh (gap-1, 2, 3, 4 không stall), các hazard dữ liệu phổ biến nhất được xử lý miễn phí, góp phần duy trì throughput cao trong hầu hết các tình huống thực tế.


\chapter{KẾT LUẬN}
\label{Chapter6}

\section{Tóm Tắt Kết Quả Đạt Được}

Khóa luận đã thiết kế và kiểm chứng thành công một hệ thống SoC RISC-V hoàn chỉnh theo đặc tả RV32I+Zicsr với pipeline 7 tầng. Các mục tiêu đề ra ban đầu đều được thực hiện đầy đủ:

\textbf{1. Lõi CPU pipeline 7 tầng:} CPU RV32I+Zicsr với pipeline IF1--IF2--ID--EX--MEM1--MEM2--WB hoạt động đúng với tất cả 47 lệnh RV32I và sáu lệnh CSR. Forwarding unit xử lý RAW hazard gap-1 đến gap-4 mà không cần stall. Hazard unit xử lý chính xác load-use (1 stall), CSR-use (1--3 stall), branch/jump flush (2 cycle), và bus stall (N cycle).

\textbf{2. Tích hợp bus chuẩn công nghiệp:} Hệ thống kết nối ngoại vi thông qua cả hai giao thức AMBA -- AXI4-Lite (1\,GHz, đồng bộ) và AHB-Lite (500\,MHz, bất đồng bộ qua CDC). Ba slave mỗi bus, mỗi slave triển khai SFR Standard Register Map 9 thanh ghi.

\textbf{3. CDC an toàn:} FIFO bất đồng bộ độ sâu 2 với con trỏ Gray code truyền dữ liệu an toàn giữa miền 1\,GHz và 500\,MHz. IRQ từ AHB domain được đồng bộ qua module 2-FF synchronizer riêng biệt.

\textbf{4. PLIC và Zicsr:} PLIC 6 nguồn tương thích chuẩn SiFive với priority, threshold và claim/complete. Zicsr xử lý đầy đủ ngắt ngoại lệ M-mode với vectored interrupt mode và precise exception semantics.

\textbf{5. Kiểm chứng toàn diện:} Hơn 630 test case tự kiểm tra cộng với 37 chương trình hệ thống, tất cả PASS. Bộ kiểm thử tuân thủ chính thức RV32I đạt 37/38 PASS (1 SKIP do giới hạn phần cứng của bài kiểm thử jal-01).

\textbf{6. SFR Standard Register Map:} Chuẩn hóa giao diện ngoại vi theo mô hình OpenTitan-inspired cho phép bất kỳ ngoại vi tuân thủ chuẩn đều có thể kết nối vào SoC mà không cần sửa RTL lõi.

\section{Hạn Chế}

Dù đạt được các mục tiêu đề ra, thiết kế còn một số hạn chế cần thừa nhận:

\begin{itemize}
    \item \textbf{Không có branch prediction:} Pipeline phải flush 2 chu kỳ mỗi khi branch taken. Với workload có nhiều branch (như vòng lặp), IPC thực tế thấp hơn lý thuyết. Thêm static prediction (predict not-taken) hoặc dynamic prediction sẽ cải thiện đáng kể.
    \item \textbf{IMEM/DMEM không có cache:} Mọi lệnh và dữ liệu đều truy cập trực tiếp SRAM 1 chu kỳ. Trong hệ thống thực, bộ nhớ ngoài có latency cao hơn nhiều, đòi hỏi cache để duy trì throughput.
    \item \textbf{Chưa tổng hợp vật lý:} Thiết kế chỉ được kiểm chứng ở mức RTL mô phỏng. Tần số 1\,GHz là mục tiêu thiết kế; việc đạt được tần số này trên công nghệ ASIC cụ thể phụ thuộc vào kết quả timing analysis sau synthesis và place-and-route.
    \item \textbf{Bộ lệnh giới hạn ở RV32I:} Không hỗ trợ nhân/chia (M extension), nguyên tử (A extension), hay dấu phẩy động (F/D). Đây là giới hạn có chủ ý để giữ phạm vi phù hợp với khóa luận cử nhân.
    \item \textbf{jal-01 không kiểm thử được qua framework:} Bài kiểm thử jal-01 trong riscv-arch-test yêu cầu IMEM $>$1.7\,MB -- không khả thi cho SoC nhúng. Tuy nhiên, lệnh JAL đã được kiểm chứng đầy đủ qua 37 chương trình system test.
\end{itemize}

\section{Hướng Phát Triển}

Dựa trên nền tảng đã xây dựng, có một số hướng phát triển tiếp theo có tiềm năng cao:

\begin{itemize}
    \item \textbf{Tổng hợp vật lý và timing closure:} Sử dụng Yosys + OpenSTA hoặc công cụ thương mại (Synopsys DC, Cadence Genus) để tổng hợp design, phân tích timing và đảm bảo closure ở 1\,GHz trên công nghệ 130nm hoặc 28nm.
    \item \textbf{Thêm branch predictor:} Triển khai BTB (Branch Target Buffer) và BHT (Branch History Table) để giảm branch penalty từ 2 xuống 0--1 cycle trong trường hợp predict đúng.
    \item \textbf{Phần mở rộng M (Multiply/Divide):} Thêm khối nhân/chia tích hợp tại tầng EX hoặc WB để hỗ trợ RV32IM -- cần thiết cho các thuật toán số học phức tạp.
    \item \textbf{Cache L1 I\$ và D\$:} Thêm instruction cache và data cache để kết nối với bộ nhớ ngoài DRAM mà không hy sinh throughput pipeline.
    \item \textbf{Hỗ trợ FPGA:} Đưa design lên board phát triển (Arty A7, Nexys A7) để kiểm chứng phần cứng thực, bao gồm tích hợp với IP ngoại vi thực tế như UART, SPI, GPIO trên chip FPGA.
    \item \textbf{Formal verification:} Áp dụng công cụ verification hình thức (như SymbiYosys) để chứng minh các bất biến pipeline (invariants) như precise exception luôn đúng, không bao giờ có data corruption.
\end{itemize}


% Công trình của tác giả (nếu không có thì comment 02 dòng dưới)
\cleardoublepage
\newpage
\phantomsection
\addcontentsline{toc}{chapter}{DANH MỤC CÔNG TRÌNH CỦA TÁC GIẢ}
\chapter*{Danh mục công trình của tác giả}
\label{Appendix1}

\begin{enumerate}
\item Tên tác giả, (năm), "Tên bài báo/tạp chí", \textit{Tên hội nghị}.
\end{enumerate}

% In tài liệu tham khảo
\cleardoublepage
\newpage
\phantomsection
\addcontentsline{toc}{chapter}{TÀI LIỆU THAM KHẢO}
\printbibheading[title={TÀI LIỆU THAM KHẢO}]

% \printbibliography[heading=subbibliography, title={Tiếng Việt}, keyword=Viet, resetnumbers=true]

\DeclareNameAlias{sortname}{last-first}
\DeclareNameAlias{default}{last-first}

\printbibliography[heading=subbibliography, title={Tiếng Anh}, notkeyword=Viet, resetnumbers=1] 
% ===================================================================== %
% CHÚ Ý: phải gán lại resetnumbers=số tài liệu tham khảo tiếng Việt + 1 %
% ===================================================================== %

% Phần phụ lục
% \include{Appendix/appendix1}
% \include{Appendix/appendix2}

\end{document} 